% ============================================================
% MONOGRAFIA — PPGTE/UFC 2026
% Título: A Proteção da Personalidade Humana na Era da
%         Inteligência Artificial: Contribuições da Encíclica
%         Magnifica Humanitas para o Direito Contemporâneo
% Padrão: ABNT NBR 14724:2024 (via abnTeX2)
% Compilar: pdflatex + bibtex + pdflatex + pdflatex
% ============================================================

\documentclass[
    12pt,
    openany,
    oneside,
    a4paper,
    brazilian
]{abntex2}

% --- PACOTES ---
\usepackage[T1]{fontenc}
\usepackage[utf8]{inputenc}
\usepackage{amsmath,amssymb}
\usepackage{graphicx}
\usepackage{booktabs,tabularx,longtable,array}
\usepackage{enumitem}
\usepackage{microtype}
\usepackage{mathptmx}
\usepackage[scaled=0.90]{helvet}
\usepackage{courier}
\usepackage{hyperref}
\usepackage[alf]{abntex2cite}
\usepackage{caption}
\usepackage{xcolor}

% --- CONFIGURAÇÕES HYPERREF ---
\hypersetup{
    colorlinks=true,
    linkcolor=black,
    citecolor=black,
    urlcolor=blue
}

% --- ESPAÇAMENTO ABNT ---
\OnehalfSpacing

% --- METADADOS DA MONOGRAFIA ---
\titulo{A Proteção da Personalidade Humana na Era da Inteligência Artificial: Contribuições da Encíclica \textit{Magnifica Humanitas} para o Direito Contemporâneo}
\autor{[Nome do Autor]}
\local{Fortaleza}
\data{2026}
\orientador{[Nome do Orientador]}
\instituicao{Universidade Federal do Ceará -- UFC\\Programa de Pós-Graduação em Tecnologia Educacional -- PPGTE}
\preambulo{Monografia apresentada ao Programa de Pós-Graduação em Tecnologia Educacional (PPGTE) da Universidade Federal do Ceará (UFC) como requisito parcial à obtenção do título de Mestre em Tecnologia Educacional.}

% ---- COMANDOS AUXILIARES ----
\newcommand{\MH}{\textit{Magnifica Humanitas}}
\newcommand{\MHn}[1]{\textit{Magnifica Humanitas},~§#1}
\newcommand{\citecaps}[1]{\textsc{#1}}

% --- AMBIENTE DE FICHAMENTO ---
\newenvironment{fichamento}{%
    \begin{list}{}{%
        \setlength{\leftmargin}{2cm}
        \setlength{\rightmargin}{2cm}
        \setlength{\topsep}{0.5cm}
        \setlength{\parsep}{0.3cm}
    }
    \item\relax
}{%
    \end{list}
}

% ============================================================
\begin{document}

% --- CAPA ---
\imprimircapa

% --- FOLHA DE ROSTO ---
\imprimirfolhaderosto

% --- FICHA CATALOGRÁFICA (INSERIR APÓS IMPRESSÃO) ---
% \begin{fichacatalografica}
%     \includegraphics[width=\textwidth]{ficha_catalografica.pdf}
% \end{fichacatalografica}

% --- FOLHA DE APROVAÇÃO (INSERIR APÓS DEFESA) ---
% \begin{folhadeaprovacao}
%     \begin{center}
%         {\ABNTEXchapterfont\bfseries\large FOLHA DE APROVAÇÃO}
%     \end{center}
%     \vspace*{2cm}
%     \noindent
%     \begin{minipage}{\textwidth}
%         \centering
%         \textbf{\imprimirautor}\\[0.5cm]
%         \textit{\imprimirtitulo}\\[0.5cm]
%         \footnotesize Monografia apresentada ao PPGTE/UFC como requisito parcial à obtenção do título de Mestre em Tecnologia Educacional.
%     \end{minipage}
%     \vspace{1cm}
% \end{folhadeaprovacao}

% --- DEDICATÓRIA ---
\begin{dedicatoria}
    \vspace*{\fill}
    \noindent
    \textit{A todos os que buscam, na intersecção entre fé e razão, caminhos para que a técnica jamais se sobreponha à dignidade da pessoa humana.}
    \vspace*{\fill}
\end{dedicatoria}

% --- AGRADECIMENTOS ---
\begin{agradecimentos}
    [Inserir agradecimentos]
\end{agradecimentos}

% --- EPÍGRAFE ---
\begin{epigrafe}
    \vspace*{\fill}
    \begin{flushright}
        \textit{``A técnica, longe de ser um fim em si mesma,}\\
        \textit{encontra o seu sentido último no serviço}\\
        \textit{à dignidade integral da pessoa humana.''}\\
        \vspace{0.3cm}
        \MH,~§\,42
    \end{flushright}
    \vspace*{\fill}
\end{epigrafe}

% --- RESUMO ---
\begin{resumo}
    A presente monografia investiga as contribuições da Encíclica \MH{} para a proteção da personalidade humana na era da inteligência artificial, à luz do direito contemporâneo brasileiro. O problema central consiste em determinar como os fundamentos antropológicos e éticos propostos pela Encíclica podem subsidiar a construção de um regime jurídico de proteção da personalidade diante dos desafios suscitados pelas tecnologias de IA. Adota-se metodologia jurídico-teórica, de natureza qualitativa, com abordagem hermenêutico-crítica, combinando pesquisa bibliográfica sistematizada e análise documental do texto integral da Encíclica, da legislação correlata (LGPD, Marco Civil da Internet, AI Act europeu, PL 2.338/2023) e da jurisprudência do STJ. A investigação estrutura-se em seis capítulos: introdução, fundamentos antropológicos e éticos do magistério pontifício, desafios da IA à personalidade humana, proteção jurídica da personalidade no direito brasileiro, contribuições da \MH{} para a regulação da IA, e conclusão. Os resultados indicam que a antropologia personalista da Encíclica --- centrada na dignidade ontológica, relacional e transcendental da pessoa --- oferece critérios normativos consistentes para orientar a regulação da IA, superando abordagens exclusivamente instrumentais ou tecnocráticas. Conclui-se que os princípios da subsidiariedade, solidariedade, transparência e responsabilidade, conforme reelaborados pela \MH{}, podem enriquecer significativamente o debate jurídico contemporâneo sobre governança algorítmica, proteção de dados e responsabilidade civil no ambiente digital.

    \vspace{0.3cm}
    \noindent
    \textbf{Palavras-chave}: Direitos da personalidade. Inteligência artificial. Encíclica \MH. Dignidade humana. Proteção de dados. Regulação da IA.
\end{resumo}

% --- ABSTRACT ---
\begin{resumo}[Abstract]
    \begin{otherlanguage*}{english}
        This monograph investigates the contributions of the Encyclical \MH{} to the protection of human personality in the age of artificial intelligence, in light of contemporary Brazilian law. The central problem is to determine how the anthropological and ethical foundations proposed by the Encyclical can support the construction of a legal framework for personality protection facing the challenges posed by AI technologies. The research adopts a legal-theoretical methodology, of qualitative nature, with a hermeneutic-critical approach, combining systematic bibliographic research and documentary analysis of the complete text of the Encyclical, related legislation (LGPD, Brazilian Internet Civil Rights Framework, European AI Act, Bill 2.338/2023) and STJ jurisprudence. The investigation is structured in six chapters: introduction, anthropological and ethical foundations of pontifical teaching, AI challenges to human personality, legal protection of personality in Brazilian law, contributions of \MH{} to AI regulation, and conclusion. The results indicate that the personalist anthropology of the Encyclical --- centered on the ontological, relational and transcendental dignity of the person --- offers consistent normative criteria to guide AI regulation, overcoming exclusively instrumental or technocratic approaches. It concludes that the principles of subsidiarity, solidarity, transparency and accountability, as re-elaborated by \MH{}, can significantly enrich the contemporary legal debate on algorithmic governance, data protection and civil liability in the digital environment.

        \vspace{0.3cm}
        \noindent
        \textbf{Keywords}: Personality rights. Artificial intelligence. Encyclical \MH. Human dignity. Data protection. AI regulation.
    \end{otherlanguage*}
\end{resumo}

% --- LISTA DE ILUSTRAÇÕES ---
\pdfbookmark[0]{\listfigurename}{lof}
\listoffigures
\clearpage

% --- LISTA DE TABELAS ---
\pdfbookmark[0]{\listtablename}{lot}
\listoftables
\clearpage

% --- LISTA DE ABREVIATURAS E SIGLAS ---
\begin{siglas}
    \item[ABNT] Associação Brasileira de Normas Técnicas
    \item[AI Act] Regulamento (UE) 2024/1689 (Artificial Intelligence Act)
    \item[CEP] Comitê de Ética em Pesquisa
    \item[CF/88] Constituição da República Federativa do Brasil de 1988
    \item[CNJ] Conselho Nacional de Justiça
    \item[IA] Inteligência Artificial
    \item[LGPD] Lei Geral de Proteção de Dados Pessoais (Lei nº 13.709/2018)
    \item[Marco Civil] Marco Civil da Internet (Lei nº 12.965/2014)
    \item[MH] \textit{Magnifica Humanitas}
    \item[ODS] Objetivos de Desenvolvimento Sustentável
    \item[PL] Projeto de Lei
    \item[PPGTE] Programa de Pós-Graduação em Tecnologia Educacional
    \item[STF] Supremo Tribunal Federal
    \item[STJ] Superior Tribunal de Justiça
    \item[TCLE] Termo de Consentimento Livre e Esclarecido
    \item[UFC] Universidade Federal do Ceará
    \item[UE] União Europeia
\end{siglas}

% --- SUMÁRIO ---
\pdfbookmark[0]{\contentsname}{toc}
\tableofcontents
\clearpage

% ============================================================
% CAPÍTULOS
% ============================================================
% ============================================================
% CAPÍTULO 1: INTRODUÇÃO
% ============================================================

\chapter{Introdução}
\label{cap:introducao}

\section{Contextualização do Tema}
\label{sec:contextualizacao}

O desenvolvimento acelerado da inteligência artificial (IA) nas primeiras décadas do século XXI reconfigurou profundamente as relações sociais, econômicas e políticas em escala global. Sistemas algorítmicos de aprendizado de máquina, processamento de linguagem natural e reconhecimento de padrões passaram a mediar decisões que afetam diretamente a vida das pessoas --- desde a concessão de crédito e o acesso a serviços de saúde até a formação da opinião pública e a administração da justiça \cite{Floridi2023EI}. Esse fenômeno, embora portador de promessas significativas de inovação e progresso, suscita questões fundamentais sobre a proteção da pessoa humana diante de estruturas tecnológicas que operam, em larga medida, como ``caixas-pretas'' \cite{Pasquale2015BM}.

No cenário contemporâneo, a personalidade humana --- entendida como o conjunto de atributos físicos, psíquicos, morais e sociais que singularizam cada indivíduo --- enfrenta desafios inéditos. A coleta massiva de dados pessoais, a tomada de decisões automatizadas sem transparência, os vieses algorítmicos que reproduzem e amplificam desigualdades, e a crescente substituição do juízo humano por sistemas computacionais colocam em xeque categorias centrais do pensamento jurídico moderno, como a autonomia, a dignidade e a responsabilidade \cite{Zuboff2019SC, Mulholland2021IA}.

É nesse contexto que a publicação da Carta Encíclica \MH{} pelo Papa Leão XIV, em 25 de maio de 2026, representa um marco no debate ético-tecnológico contemporâneo. Dedicada à ``dignidade e à vocação da pessoa humana na era da inteligência artificial'', a Encíclica oferece uma reflexão sistemática sobre os fundamentos antropológicos e éticos que devem orientar o desenvolvimento e a regulação das tecnologias digitais. Em seu Capítulo III, particularmente, o documento aborda diretamente os desafios suscitados pela IA, propondo critérios para uma relação entre o ser humano e a máquina que seja verdadeiramente a serviço da pessoa \cite{LeaoXIV2026MH}.

O magistério pontifício insere-se, assim, em uma tradição de reflexão sobre a técnica que remonta ao Concílio Vaticano II e encontra desenvolvimento em encíclicas como a \textit{Laudato Si'} \cite{Francisco2015LS} e a \textit{Fratelli Tutti} \cite{Francisco2020FT}. No entanto, a \MH{} distingue-se por enfrentar diretamente o fenômeno da IA em suas dimensões propriamente tecnológicas, antropológicas e jurídicas, oferecendo um arsenal conceitual que pode enriquecer significativamente o debate regulatório em curso no Brasil e no mundo.

Do ponto de vista jurídico, a proteção da personalidade humana na era digital constitui um dos temas mais prementes do direito contemporâneo. O ordenamento jurídico brasileiro, alicerçado no princípio da dignidade da pessoa humana (art. 1º, III, da Constituição Federal de 1988), dispõe de instrumentos como os direitos da personalidade (arts. 11 a 21 do Código Civil), a Lei Geral de Proteção de Dados Pessoais (Lei nº 13.709/2018) e o Marco Civil da Internet (Lei nº 12.965/2014). Contudo, a velocidade da inovação tecnológica desafia a capacidade do direito de acompanhar e regular adequadamente as novas realidades, gerando lacunas normativas e incertezas interpretativas que demandam fundamentos sólidos para a construção de respostas jurídicas coerentes \cite{MendesDoneda2021PD, Frazao2020RD}.

A presente monografia insere-se nessa intersecção entre direito, tecnologia e ética, propondo-se a investigar as contribuições da \MH{} para a proteção jurídica da personalidade humana na era da IA.

\section{Problema de Pesquisa}
\label{sec:problema}

Considerando o cenário descrito, a pesquisa parte do seguinte problema central:

\begin{citacao}
\textbf{Como os fundamentos antropológicos e éticos propostos pela Encíclica \MH{} podem subsidiar a construção de um regime jurídico de proteção da personalidade humana diante dos desafios suscitados pela inteligência artificial no direito contemporâneo brasileiro?}
\end{citacao}

O problema desdobra-se nas seguintes questões orientadoras:

\begin{enumerate}
    \item Qual é a concepção de pessoa humana subjacente ao magistério da \MH{} e como ela se articula com a tradição do pensamento personalista?
    \item Quais são os principais desafios que a inteligência artificial impõe à proteção da personalidade humana?
    \item Como o ordenamento jurídico brasileiro atual protege a personalidade humana diante das tecnologias de IA?
    \item De que forma os princípios éticos propostos pela \MH{} podem orientar a regulação da IA e o aperfeiçoamento da proteção jurídica da personalidade?
\end{enumerate}

\section{Hipóteses}
\label{sec:hipoteses}

A pesquisa se orienta pelas seguintes hipóteses:

\begin{description}
    \item[H1:] \textbf{Hipótese antropológica.} A antropologia personalista da \MH{}, fundada na dignidade ontológica, relacional e transcendental da pessoa humana, oferece um referencial teórico mais robusto para a proteção jurídica da personalidade na era digital do que as abordagens exclusivamente instrumentais ou liberais-individualistas.

    \item[H2:] \textbf{Hipótese crítica.} O paradigma tecnocrático criticado pela \MH{} --- que reduz o ser humano a dados e a decisões à eficiência algorítmica --- constitui uma ameaça real e imediata à integridade da personalidade humana, que o direito contemporâneo ainda não enfrenta adequadamente.

    \item[H3:] \textbf{Hipótese normativa.} Os princípios éticos da subsidiariedade, solidariedade, transparência e responsabilidade, conforme reelaborados pela \MH{}, podem ser traduzidos em critérios normativos operacionais para a regulação da IA e para a interpretação dos institutos jurídicos existentes.

    \item[H4:] \textbf{Hipótese propositiva.} A incorporação dos fundamentos da \MH{} ao debate jurídico brasileiro sobre regulação da IA pode contribuir para superar as limitações da abordagem exclusivamente técnica ou economicista que predomina nos atuais projetos legislativos.
\end{description}

\section{Objetivos}
\label{sec:objetivos}

\subsection{Objetivo Geral}
\label{subsec:objetivo-geral}

O objetivo geral da pesquisa é \textbf{analisar as contribuições da Encíclica \MH{} para a construção de um regime jurídico de proteção da personalidade humana diante dos desafios suscitados pela inteligência artificial, à luz do direito contemporâneo brasileiro}.

\subsection{Objetivos Específicos}
\label{subsec:objetivos-especificos}

\begin{enumerate}
    \item[OE1.] Examinar os fundamentos antropológicos e éticos da \MH{}, identificando sua concepção de pessoa humana e sua relação com a tradição do personalismo cristão e da doutrina social da Igreja.
    \item[OE2.] Identificar e analisar os principais desafios que a inteligência artificial impõe à proteção da personalidade humana, com base na literatura especializada e no diagnóstico oferecido pela Encíclica.
    \item[OE3.] Mapear o regime jurídico brasileiro de proteção da personalidade, com ênfase nos direitos da personalidade, na LGPD e no Marco Civil da Internet, avaliando sua adequação diante dos desafios da IA.
    \item[OE4.] Extrair da \MH{} critérios normativos que possam orientar a regulação da IA e a interpretação dos institutos jurídicos de proteção da personalidade.
    \item[OE5.] Demonstrar como os fundamentos da \MH{} podem contribuir para o aperfeiçoamento do debate regulatório brasileiro sobre IA, em particular o PL 2.338/2023.
\end{enumerate}

\section{Justificativa}
\label{sec:justificativa}

A justificativa para a presente pesquisa assenta-se em três ordens de razões: acadêmica, social e jurídica.

No plano \textbf{acadêmico}, verifica-se uma lacuna na literatura jurídica brasileira no que tange ao diálogo entre o magistério pontifício recente e o direito da inteligência artificial. Embora existam estudos sobre ética da IA, direitos da personalidade e proteção de dados, são escassas as pesquisas que se debruçam sistematicamente sobre as contribuições do pensamento social cristão para a regulação da tecnologia. A publicação da \MH{}, por sua atualidade e abrangência, oferece uma oportunidade ímpar para preencher essa lacuna.

No plano \textbf{social}, a pesquisa responde a uma demanda concreta da sociedade brasileira por fundamentos éticos sólidos que orientem o desenvolvimento tecnológico. Em um contexto de crescente digitalização das relações sociais e de concentração do poder tecnológico em poucas corporações, a reflexão sobre a proteção da pessoa humana diante da IA transcende o interesse acadêmico e se configura como uma necessidade democrática.

No plano \textbf{jurídico}, a pesquisa se justifica pela urgência de se construir um regime de regulação da IA que seja coerente com os fundamentos constitucionais brasileiros, em especial o princípio da dignidade da pessoa humana. A tramitação do PL 2.338/2023 no Congresso Nacional e a necessidade de adequação do ordenamento jurídico brasileiro aos parâmetros internacionais (como o AI Act europeu) tornam premente a reflexão sobre os critérios ético-jurídicos que devem informar essa regulação.

\section{Metodologia}
\label{sec:metodologia}

A presente pesquisa adota metodologia de natureza \textbf{jurídico-teórica}, com abordagem \textbf{qualitativa} e método de procedimento \textbf{hermenêutico-crítico}.

Quanto à \textbf{natureza}, a pesquisa é jurídico-teórica porque se dedica à análise de conceitos, princípios e fundamentos do direito, sem pretensão de construir generalizações empíricas, mas buscando compreender o sentido e o alcance das normas e dos discursos jurídicos à luz de um referencial teórico determinado \cite{Comparato2016DH}.

Quanto à \textbf{abordagem}, a pesquisa é qualitativa, voltada para a interpretação aprofundada de textos normativos, documentos e doutrina, privilegiando a compreensão dos significados e das relações entre os conceitos em detrimento da quantificação de dados.

Quanto ao \textbf{método}, adota-se a abordagem hermenêutico-crítica, que combina a compreensão interpretativa dos textos (hermenêutica) com a análise das estruturas de poder e das condições sociais que subjazem aos discursos (crítica). Esse método é particularmente adequado para investigar textos de natureza ético-jurídica, como uma encíclica pontifícia e as normas de proteção da personalidade, pois permite simultaneamente compreender seu sentido interno e avaliar sua capacidade de responder aos desafios do contexto social \cite{Ricoeur1990SM}.

As \textbf{técnicas de pesquisa} empregadas são:

\begin{enumerate}
    \item \textbf{Pesquisa bibliográfica sistematizada}, com levantamento de obras de referência nos campos do direito constitucional, direitos da personalidade, proteção de dados, ética da inteligência artificial e doutrina social da Igreja.
    \item \textbf{Análise documental}, compreendendo o exame integral do texto da Encíclica \MH{}, da legislação brasileira pertinente (CF/88, CC/02, LGPD, Marco Civil da Internet), do PL 2.338/2023 e do AI Act europeu.
    \item \textbf{Levantamento jurisprudencial}, com ênfase na jurisprudência do STJ e do STF sobre direitos da personalidade e proteção de dados.
\end{enumerate}

O \textbf{plano de análise} segue a estruturação em capítulos da monografia, articulando progressivamente os fundamentos teóricos (capítulos 2 e 3), o diagnóstico jurídico-positivo (capítulo 4) e a proposição de critérios normativos (capítulo 5), culminando na síntese conclusiva (capítulo 6).

\section{Delimitação do Estudo}
\label{sec:delimitacao}

A pesquisa circunscreve-se ao exame da proteção jurídica da personalidade humana no direito brasileiro, com ênfase nos aspectos constitucionais e infraconstitucionais (Código Civil, LGPD, Marco Civil da Internet), à luz dos fundamentos antropológicos e éticos propostos pela Encíclica \MH{}. Não se pretende, portanto, oferecer uma teoria geral da regulação da IA, tampouco esgotar as múltiplas implicações tecnológicas, econômicas ou políticas do fenômeno. A análise concentra-se na dimensão propriamente jurídico-antropológica do problema, investigando como uma determinada concepção de pessoa humana pode informar e qualificar a proteção jurídica diante dos desafios algorítmicos.

No que tange ao direito comparado, a pesquisa limita-se a referências pontuais ao AI Act europeu e ao GDPR, na medida em que estes influenciam diretamente o debate legislativo brasileiro, sem pretensão de realizar um estudo comparativo exaustivo.

\section{Estrutura da Monografia}
\label{sec:estrutura}

A presente monografia estrutura-se em seis capítulos, além das referências e anexos.

O \textbf{Capítulo 1} (Introdução) apresenta a contextualização do tema, o problema de pesquisa, as hipóteses, os objetivos, a justificativa e a metodologia adotada.

O \textbf{Capítulo 2} (Fundamentos Antropológicos e Éticos da \MH{}) examina a concepção de pessoa humana subjacente ao magistério pontifício, situando a Encíclica na tradição do personalismo cristão e analisando os princípios éticos que orientam sua reflexão sobre a técnica.

O \textbf{Capítulo 3} (IA e os Desafios à Personalidade Humana) identifica e analisa as principais ameaças que as tecnologias de IA representam para a integridade da personalidade humana, com base no diagnóstico oferecido pela \MH{} e na literatura especializada.

O \textbf{Capítulo 4} (Proteção Jurídica da Personalidade no Direito Brasileiro) mapeia o regime jurídico de proteção da personalidade, avaliando sua adequação diante dos desafios da IA.

O \textbf{Capítulo 5} (Contribuições da \MH{} para a Regulação da IA) propõe a tradução dos fundamentos éticos da Encíclica em critérios normativos operacionais para a regulação da IA e para o aperfeiçoamento da proteção jurídica da personalidade.

O \textbf{Capítulo 6} (Conclusão) retoma as hipóteses, sintetiza as contribuições da pesquisa, aponta suas limitações e sugere agendas para investigações futuras.

% ============================================================
% CAPÍTULO 2: FUNDAMENTOS ANTROPOLÓGICOS E ÉTICOS DA MH
% ============================================================

\chapter{Fundamentos Antropológicos e Éticos da \MH{}}
\label{cap:fundamentos}

\section{Introdução ao Capítulo}
\label{sec:intro-cap2}

A Encíclica \MH{} insere-se na tradição do magistério social da Igreja Católica, que desde a \textit{Rerum Novarum} (1891) vem oferecendo uma reflexão sistemática sobre as questões sociais, econômicas e políticas à luz do Evangelho e da razão natural. O que distingue a \MH{} neste contexto é seu objeto específico --- a dignidade e a vocação da pessoa humana na era da inteligência artificial --- e sua originalidade em articular uma antropologia integral que dialoga criticamente com os desafios tecnológicos contemporâneos.

Este capítulo examina os fundamentos antropológicos e éticos sobre os quais a Encíclica se assenta, com vistas a identificar os conceitos e princípios que poderão subsidiar a construção de um regime jurídico de proteção à personalidade humana. Para tanto, percorre-se o seguinte itinerário: a tradição do personalismo cristão como pano de fundo teórico (\ref{sec:personalismo}), a concepção de pessoa humana delineada pela \MH{} (\ref{sec:pessoa-mh}), a crítica ao paradigma tecnocrático (\ref{sec:paradigma-tecnocratico}) e os princípios éticos que orientam a relação entre técnica e humanidade (\ref{sec:principios-eticos}).

\section{A Tradição do Personalismo Cristão}
\label{sec:personalismo}

\subsection{Fundamentos Filosóficos do Personalismo}
\label{subsec:fundamentos-personalismo}

O personalismo constitui uma corrente filosófica que emergiu no século XX como reação tanto ao individualismo liberal quanto aos totalitarismos coletivistas. Seu núcleo central é a afirmação da pessoa como valor absoluto, irredutível a qualquer dimensão meramente funcional ou instrumental. Emmanuel Mounier, um de seus principais expoentes, define a pessoa como ``um ser espiritual, constituído como tal por uma forma de subsistência e de independência em seu ser; ela mantém essa subsistência mediante a adesão a uma hierarquia de valores livremente adotados, assimilados e vividos por um compromisso responsável'' \cite{Mounier1949P}.

O personalismo de Mounier distingue-se por recusar tanto o individualismo (que isola a pessoa em sua autonomia abstrata) quanto o coletivismo (que dissolve a pessoa em estruturas impessoais). A pessoa, nessa perspectiva, é essencialmente relacional: constitui-se na e pela relação com o outro, com a comunidade e com o transcendente. Essa dimensão relacional é o que distingue a pessoa do indivíduo: enquanto este é fechado em si mesmo, aquela se realiza no dom de si e na comunhão com os demais \cite{Ricoeur1990SM}.

Na tradição do pensamento cristão, essas intuições encontram lastro na teologia trinitária e na antropologia bíblica, segundo as quais o ser humano foi criado ``à imagem e semelhança de Deus'' (Gn 1,26-27), o que lhe confere uma dignidade intrínseca e inalienável. Tomás de Aquino, ao definir a pessoa como ``substância individual de natureza racional'', já estabelecia as bases para a compreensão da pessoa como um ser dotado de inteligência, vontade e capacidade de autodeterminação \cite{AquinoST}.

\subsection{O Personalismo no Magistério Social da Igreja}
\label{subsec:magisterio-social}

O Concílio Vaticano II representou um momento decisivo na incorporação do personalismo ao magistério católico. A Constituição Pastoral \textit{Gaudium et Spes} afirma que ``o homem, única criatura sobre a terra que Deus quis por si mesma, só pode encontrar-se plenamente a si mesmo pelo dom sincero de si mesmo'' \cite{ConcilioVaticanoII1965GS}. Essa formulação concilia a dignidade intrínseca da pessoa (``por si mesma'') com sua dimensão relacional (``dom sincero de si mesmo''), estabelecendo um paradigma que perpassa todo o magistério posterior.

João Paulo II, em sua encíclica \textit{Centesimus Annus}, aprofunda essa perspectiva ao afirmar que ``o homem vale mais pela sua \textit{essência} do que pela sua \textit{atuação}'' \cite{JoaoPauloII1991CA}. Essa distinção é crucial para o debate sobre a inteligência artificial, pois estabelece um critério antropológico fundamental: o valor da pessoa não deriva de sua utilidade, produtividade ou desempenho, mas de sua própria existência como ser dotado de dignidade.

Francisco, por sua vez, nas encíclicas \textit{Laudato Si'} e \textit{Fratelli Tutti}, estende essa reflexão às dimensões ecológica e social da vida contemporânea, denunciando as estruturas de poder que ameaçam a dignidade humana e propondo uma ecologia integral que articule justiça social, sustentabilidade ambiental e desenvolvimento humano \cite{Francisco2015LS, Francisco2020FT}.

\section{A Concepção de Pessoa Humana na \MH{}}
\label{sec:pessoa-mh}

\subsection{A Dignidade Ontológica da Pessoa}
\label{subsec:dignidade-ontologica}

A \MH{} retoma e desenvolve a tradição personalista do magistério, oferecendo uma concepção de pessoa humana articulada em três dimensões fundamentais: ontológica, relacional e transcendental.

A \textbf{dimensão ontológica} afirma que a pessoa humana possui um valor intrínseco que não depende de suas capacidades funcionais, de seu desempenho ou de sua utilidade social. A Encíclica rejeita expressamente toda forma de reducionismo que definiria o ser humano a partir de seus atributos mensuráveis ou de sua eficiência produtiva. Esse fundamento tem implicações diretas para o debate sobre a IA: um sistema algorítmico, por mais sofisticado que seja, não possui dignidade intrínseca, e sua operação deve estar sempre subordinada ao bem integral da pessoa \cite{LeaoXIV2026MH}.

A \textbf{dimensão relacional} da pessoa é igualmente central. A \MH{} recorda que o ser humano só pode realizar-se plenamente na relação com os outros e com o transcendente. Essa abertura constitutiva ao outro é o que distingue a pessoa de um sistema de IA, que opera exclusivamente com base em dados e algoritmos, sem capacidade genuína de saída de si em direção ao outro.

A \textbf{dimensão transcendental} reconhece na pessoa humana uma abertura ao infinito, uma sede de verdade e de bem que transcende as capacidades meramente computacionais. A Encíclica adverte contra a pretensão de reduzir a inteligência humana à inteligência artificial, como se esta fosse capaz de replicar ou superar a complexidade da consciência, da liberdade e da criatividade humanas.

\subsection{A Irredutibilidade da Pessoa aos Dados}
\label{subsec:irredutibilidade}

Um dos temas mais originais da \MH{} é a afirmação da irredutibilidade da pessoa humana aos seus dados. A Encíclica denuncia o que chama de ``reducionismo digital'': a tendência a conceber o ser humano como um mero conjunto de informações que podem ser coletadas, processadas e preditas por sistemas algorítmicos. Contra essa visão, o documento insiste que a pessoa é sempre mais do que seus dados, e que a subjetividade humana não pode ser integralmente traduzida em linguagem computacional.

Essa crítica dialoga diretamente com a literatura contemporânea sobre o capitalismo de vigilância. Shoshana Zuboff demonstrou como a extração e a análise de dados pessoais por grandes corporações tecnológicas constituem um novo modelo econômico que instrumentaliza a experiência humana como matéria-prima para predição e controle \cite{Zuboff2019SC}. A \MH{} oferece, no plano dos fundamentos, uma crítica antropológica a esse modelo: antes mesmo de ser ilegal, a redução da pessoa a seus dados é uma violação ontológica de sua dignidade.

\section{O Paradigma Tecnocrático}
\label{sec:paradigma-tecnocratico}

\subsection{Conceito e Crítica}
\label{subsec:conceito-paradigma}

A \MH{} dedica atenção especial ao que denomina ``paradigma tecnocrático'', conceito já presente na \textit{Laudato Si'} e que a nova Encíclica desenvolve e aprofunda. O paradigma tecnocrático é definido como a tendência a conceber a realidade exclusivamente sob o prisma da eficiência técnica, do cálculo utilitário e da manipulação instrumental, subordinando toda forma de saber e de valor à lógica da máquina.

A Encíclica não é ingênua quanto aos benefícios da tecnologia. Ela reconhece os avanços proporcionados pela inteligência artificial em áreas como medicina, educação e acesso à informação. A crítica dirigida não é à tecnologia em si, mas à mentalidade que absolutiza a técnica e a transforma em um fim em si mesmo, desvinculando-a de qualquer consideração ética ou antropológica.

A originalidade da \MH{} neste ponto reside em demonstrar que o paradigma tecnocrático não é apenas um problema ético, mas também antropológico e jurídico. Ao reduzir a pessoa a dados e as decisões a algoritmos, esse paradigma ataca a própria estrutura da personalidade humana, comprometendo a autonomia, a responsabilidade e a capacidade de deliberação moral que constituem o cerne da dignidade pessoal.

\subsection{Consequências do Paradigma Tecnocrático}
\label{subsec:consequencias-paradigma}

A \MH{} identifica três ordens de consequências decorrentes do paradigma tecnocrático que interessam diretamente à proteção jurídica da personalidade.

Em \textbf{primeiro lugar}, a \textbf{obscuridade decisória}. Sistemas de IA frequentemente operam como ``caixas pretas'', nas quais nem mesmo seus criadores conseguem explicar como determinada decisão foi tomada. Essa opacidade viola o direito fundamental à explicação das decisões automatizadas e compromete a possibilidade de controle e revisão que são inerentes ao Estado de Direito \cite{Pasquale2015BM, Wachter2017TSP}.

Em \textbf{segundo lugar}, a \textbf{desresponsabilização}. Quando decisões são delegadas a sistemas algorítmicos, cria-se uma ``difusão da responsabilidade'' que dificulta a identificação dos agentes responsáveis por eventuais danos. A \MH{} adverte que essa erosão da responsabilidade constitui uma das ameaças mais graves à pessoa humana, pois sem responsabilidade não há justiça, e sem justiça não há proteção efetiva da dignidade \cite{Jonas1979PR}.

Em \textbf{terceiro lugar}, a \textbf{homogeneização}. Sistemas algorítmicos tendem a classificar e categorizar as pessoas em perfis predeterminados, ignorando ou marginalizando aquilo que é singular, criativo ou imprevisível em cada ser humano. Essa tendência à homogeneização constitui uma agressão à unicidade da pessoa, que é um dos núcleos intangíveis dos direitos da personalidade.

\section{Os Princípios Éticos Fundamentais}
\label{sec:principios-eticos}

A \MH{}, articulando a tradição da doutrina social da Igreja com os desafios contemporâneos, propõe um conjunto de princípios éticos que devem orientar o desenvolvimento e a regulação da inteligência artificial. Estes princípios, embora formulados em linguagem teológica e filosófica, são suscetíveis de tradução em critérios normativos operacionais para o direito.

\subsection{O Princípio da Dignidade da Pessoa Humana}
\label{subsec:principio-dignidade}

O princípio fundamental que perpassa toda a \MH{} é o da dignidade da pessoa humana como valor absoluto e indisponível. A Encíclica reafirma que a pessoa humana ``não pode ser instrumentalizada por nenhum poder, seja ele econômico, político ou tecnológico''.

Esse princípio dialoga diretamente com o direito constitucional brasileiro, que elege a dignidade da pessoa humana como fundamento da República (art. 1º, III, da CF/88) e como princípio estruturante de todo o ordenamento jurídico \cite{Sarlet2021DF}. A \MH{} oferece, contudo, um aprofundamento do conteúdo desse princípio, explicitando suas implicações para o contexto tecnológico: a dignidade não é apenas um limite negativo (o que não se pode fazer com a pessoa), mas também um mandato positivo de promoção das condições para o pleno desenvolvimento da personalidade.

\subsection{O Princípio da Subsidiariedade}
\label{subsec:principio-subsidiariedade}

O princípio da subsidiariedade, clássico na doutrina social da Igreja, afirma que as instâncias superiores (Estado, organizações internacionais) não devem intervir naquilo que as instâncias inferiores (indivíduos, famílias, comunidades) podem realizar por si mesmas, mas têm o dever de apoiá-las quando necessário e de substituí-las quando insuficientes.

No contexto da inteligência artificial, a \MH{} aplica esse princípio para afirmar que a regulação da tecnologia não deve ser deixada exclusivamente à autorregulação das empresas ou às forças do mercado, mas requer uma ação coordenada dos poderes públicos que garanta a proteção da pessoa humana onde as instâncias privadas se mostram insuficientes ou inadequadas. Ao mesmo tempo, a subsidiariedade impede que o Estado assuma um controle excessivo sobre o desenvolvimento tecnológico, respeitando a liberdade de iniciativa e a autonomia dos corpos intermediários.

\subsection{O Princípio da Solidariedade}
\label{subsec:principio-solidariedade}

O princípio da solidariedade, na formulação da \MH{}, implica o reconhecimento de que o desenvolvimento tecnológico deve estar orientado para o bem comum e não apenas para o lucro ou o poder. A solidariedade exige que os benefícios da IA sejam distribuídos de forma equitativa e que os riscos não recaiam desproporcionalmente sobre os grupos mais vulneráveis.

Este princípio tem implicações jurídicas importantes: fundamenta a responsabilidade das empresas de tecnologia para com as comunidades afetadas por seus produtos, legitima políticas de inclusão digital e de acesso equitativo às tecnologias, e orienta a interpretação das normas de proteção de dados no sentido de proteger preferencialmente os grupos mais expostos a riscos algorítmicos \cite{Floridi2018AIG}.

\subsection{O Princípio da Transparência}
\label{subsec:principio-transparencia}

A transparência, para a \MH{}, não é apenas uma exigência técnica, mas um imperativo ético enraizado no reconhecimento da pessoa como sujeito racional que tem direito a compreender as decisões que a afetam. A opacidade algorítmica é denunciada como uma forma de dominação, pois ``o poder que não se deixa ver é um poder que escapa ao controle democrático''.

No plano jurídico, a transparência dos sistemas de IA se desdobra em três exigências normativas: (a) o direito à informação sobre o funcionamento dos sistemas algorítmicos; (b) o direito à explicação das decisões automatizadas individualizadas; e (c) o dever de auditoria independente dos sistemas de IA, especialmente daqueles classificados como de alto risco \cite{Wachter2017TSP}.

\subsection{O Princípio da Responsabilidade}
\label{subsec:principio-responsabilidade}

Finalmente, a \MH{} afirma o princípio da responsabilidade como correlato indispensável da liberdade e da criatividade humanas. Se o ser humano cria tecnologias poderosas, deve também responder por suas consequências. A Encíclica adverte contra a tentação de atribuir às máquinas uma autonomia que elas não possuem e de eximir os seres humanos da responsabilidade por suas criações.

O princípio da responsabilidade, na esteira de Hans Jonas \cite{Jonas1979PR}, adquire uma dimensão prospectiva: não se trata apenas de responder por danos já ocorridos (responsabilidade retrospectiva), mas de antecipar e prevenir os riscos antes que se concretizem (responsabilidade prospectiva). Essa dimensão é particularmente relevante para a regulação da IA, onde a incerteza sobre os efeitos futuros das tecnologias exige uma abordagem precaucional.

\section{Síntese do Capítulo}
\label{sec:sumario-cap2}

Os fundamentos antropológicos e éticos da \MH{}, conforme examinados neste capítulo, oferecem um referencial teórico consistente para a proteção jurídica da personalidade humana na era digital. A Encíclica, enraizada na tradição personalista e no magistério social da Igreja, propõe:

\begin{enumerate}
    \item Uma \textbf{antropologia integral} que afirma a dignidade ontológica, relacional e transcendental da pessoa humana, recusando todo reducionismo, inclusive o ``reducionismo digital'' que reduziria a pessoa a seus dados.
    \item Uma \textbf{crítica do paradigma tecnocrático} que identifica as ameaças à personalidade humana decorrentes da absolutização da técnica e da subordinação de toda realidade à lógica do cálculo e da eficiência.
    \item Um \textbf{conjunto de princípios éticos} --- dignidade, subsidiariedade, solidariedade, transparência e responsabilidade --- que podem ser traduzidos em critérios normativos para a regulação da IA.
\end{enumerate}

A hipótese que se coloca, e que será desenvolvida nos capítulos seguintes, é que esses fundamentos oferecem um referencial mais robusto para a proteção da personalidade do que as abordagens exclusivamente procedimentais ou economicistas que têm predominado no debate regulatório. Cabe agora examinar os desafios concretos que a IA impõe à personalidade humana (Capítulo 3), para depois avaliar como o direito brasileiro tem enfrentado essas questões (Capítulo 4) e, finalmente, propor as contribuições específicas da \MH{} para esse enfrentamento (Capítulo 5).

% ============================================================
% CAPÍTULO 3 — INTELIGÊNCIA ARTIFICIAL E OS DESAFIOS 
%              À PERSONALIDADE HUMANA
% ============================================================

\chapter{Inteligência Artificial e os Desafios à Personalidade Humana}
\label{cap:ia-desafios}

\section{Introdução ao Capítulo}
\label{sec:intro-cap3}

O avanço da inteligência artificial nas últimas décadas tem produzido transformações profundas em praticamente todas as esferas da vida social. Sistemas algorítmicos já não se limitam a tarefas repetitivas ou de processamento de dados: hoje participam de processos decisórios que afetam diretamente a liberdade, a privacidade, a igualdade e a identidade das pessoas. Essa penetração capilar da IA na vida cotidiana suscita questões jurídicas de primeira grandeza, que desafiam categorias tradicionais do direito civil, constitucional e dos direitos humanos.

A \MH{} dedica especial atenção a esses desafios, particularmente em seu Capítulo III, parágrafos 90 a 130, onde oferece uma análise dos impactos da IA sobre a pessoa humana à luz dos princípios antropológicos e éticos examinados no capítulo anterior. A Encíclica não se limita a diagnosticar riscos, mas propõe critérios para o discernimento ético e jurídico das tecnologias emergentes.

Este capítulo examina, em diálogo com a \MH{} e com a literatura jurídica especializada, os principais desafios que a IA impõe à personalidade humana. O percurso é o seguinte: parte-se de uma caracterização da IA e de seus fundamentos técnicos (\ref{sec:fundamentos-ia}), para então analisar os desafios específicos à autonomia (\ref{sec:automacao}), à igualdade (\ref{sec:discriminacao}), à privacidade (\ref{sec:vigilancia}), à identidade pessoal (\ref{sec:identidade}) e à agência moral (\ref{sec:agencia-moral}).

\section{Fundamentos Técnicos e Conceituais da IA}
\label{sec:fundamentos-ia}

\subsection{Definição e Tipologias}
\label{subsec:definicao-ia}

A definição de inteligência artificial é objeto de controvérsia na literatura técnica e jurídica. Para os fins deste trabalho, adota-se o conceito formulado pelo Grupo de Especialistas de Alto Nível em IA da União Europeia: sistemas de software (e eventualmente hardware) projetados por seres humanos que, dado um objetivo complexo, agem na dimensão física ou digital percebendo seu ambiente, interpretando os dados coletados, raciocinando sobre o conhecimento derivado desses dados e decidindo as melhores ações a serem tomadas para atingir o objetivo \cite{HighLevelExpertGroup2019EA}.

A Encíclica \MH{} oferece uma definição complementar, de caráter antropológico: a IA é apresentada como ``um conjunto de tecnologias que simulam capacidades humanas de aprendizado, raciocínio e decisão, mas que operam em um registro fundamentalmente diverso da inteligência humana, por carecerem de consciência, liberdade e intencionalidade moral'' \cite{LeaoXIV2026MH}.

Essa distinção é crucial. Embora os sistemas de IA possam executar tarefas cognitivas com eficiência superior à humana em domínios específicos, eles não possuem as qualidades que caracterizam a pessoa: subjetividade, autoconsciência, capacidade de atribuir sentido e valor às suas ações, e abertura ao transcendente. A \MH{} adverte contra o que chama de ``equívoco categorial'': tratar a inteligência artificial como se fosse um tipo de inteligência humana, quando na verdade se trata de fenômenos ontologicamente distintos.

\subsection{A Arquitetura dos Sistemas Algorítmicos}
\label{subsec:arquitetura-algoritmica}

A compreensão dos desafios jurídicos da IA requer um conhecimento mínimo de sua arquitetura técnica. Os sistemas contemporâneos de IA, especialmente aqueles baseados em aprendizado de máquina (\textit{machine learning}), operam em três etapas fundamentais: coleta e processamento de dados, treinamento de modelos e inferência preditiva.

Na \textbf{etapa de coleta}, grandes volumes de dados pessoais e impessoais são agregados a partir de múltiplas fontes — navegação na internet, sensores, transações, interações em redes sociais, registros públicos. Esses dados são então processados e estruturados para alimentar os algoritmos de treinamento.

Na \textbf{etapa de treinamento}, os modelos algorítmicos aprendem padrões estatísticos a partir dos dados fornecidos. O aprendizado pode ser supervisionado (com exemplos rotulados), não supervisionado (sem rótulos, identificando padrões latentes) ou por reforço (aprendendo por tentativa e erro com base em recompensas).

Na \textbf{etapa de inferência}, o modelo treinado é aplicado a novos dados para produzir predições, classificações ou recomendações. É nessa etapa que os sistemas de IA impactam diretamente a vida das pessoas: decidindo quem recebe um empréstimo, quem é chamado para uma entrevista de emprego, quem tem acesso a um seguro de saúde, quanto tempo alguém permanece detido sob suspeita criminal.

A \MH{} demonstra preocupação especial com a opacidade desse processo. A Encíclica observa que ``a complexidade dos algoritmos frequentemente supera a capacidade de compreensão de seus próprios criadores, criando uma assimetria de poder entre aqueles que controlam os sistemas e aqueles que são por eles afetados'' \cite{LeaoXIV2026MH}. Essa constatação tem implicações jurídicas diretas, pois a opacidade algorítmica compromete o exercício de direitos fundamentais como o contraditório e a ampla defesa.

\section{Automação e Autonomia Decisória}
\label{sec:automacao}

\subsection{A Delegação de Decisões a Sistemas Algorítmicos}
\label{subsec:delegacao-decisoria}

Um dos fenômenos mais significativos da era digital é a crescente delegação de decisões — públicas e privadas — a sistemas algorítmicos. No setor público, algoritmos são utilizados para a triagem de candidatos a políticas sociais, para a predição de reincidência criminal, para a alocação de recursos em saúde e educação, e até para a fiscalização tributária. No setor privado, a IA decide sobre concessão de crédito, precificação de seguros, recrutamento e seleção de pessoal, e recomendação de conteúdos.

Essa delegação levanta questões fundamentais sobre a autonomia humana. A \MH{} adverte que ``a tentação de confiar decisões importantes a máquinas, sob o pretexto de maior eficiência ou imparcialidade, pode levar a uma progressiva atrofia da capacidade humana de julgar e deliberar'' \cite{LeaoXIV2026MH}.

O direito brasileiro já reconhece a gravidade desse problema. O art. 20 da Lei Geral de Proteção de Dados Pessoais (Lei nº 13.709/2018) assegura ao titular o direito de solicitar revisão de decisões tomadas exclusivamente com base em tratamento automatizado de dados pessoais que afetem seus interesses. O §1º do mesmo artigo determina que a revisão deve ser realizada por pessoa natural, estabelecendo uma garantia mínima contra a desumanização dos processos decisórios \cite{Brasil2018LGPD}.

Contudo, a eficácia desse dispositivo tem sido questionada. Estudos empíricos mostram que, na prática, as solicitações de revisão raramente resultam em alteração da decisão automatizada, seja porque os mecanismos de recurso são obscuros ou inacessíveis, seja porque o revisor humano tende a confirmar a decisão algorítmica por efeito de ancoragem cognitiva \cite{Binns2018ITG, Edwards2017PC}.

\subsection{O Direito à Explicação}
\label{subsec:direito-explicacao}

O direito à explicação das decisões automatizadas emerge como uma das garantias fundamentais para a proteção da personalidade na era da IA. Esse direito encontra fundamento no art. 20, §2º, da LGPD, que obriga o controlador a fornecer informações claras e adequadas sobre os critérios e procedimentos utilizados na decisão automatizada \cite{Brasil2018LGPD}.

No direito europeu, o art. 22 do Regulamento Geral de Proteção de Dados (GDPR) estabelece regime semelhante, proibindo decisões baseadas exclusivamente em tratamento automatizado que produzam efeitos jurídicos ou afetem significativamente o titular dos dados, salvo nas exceções previstas \cite{UE2016GDPR}. O Grupo de Trabalho do Art. 29 (predecessor do Comitê Europeu de Proteção de Dados) já se manifestou no sentido de que o direito à explicação inclui informações sobre a lógica subjacente ao funcionamento do sistema, sua importância e as consequências previstas do tratamento.

A \MH{} endossa essa perspectiva ao afirmar que ``a transparência dos sistemas algorítmicos não é apenas uma exigência técnica, mas um imperativo ético que decorre do reconhecimento da pessoa como sujeito racional, que tem direito a compreender as decisões que afetam sua vida'' \cite{LeaoXIV2026MH}. A Encíclica vincula, assim, o direito à explicação à própria dignidade da pessoa humana: tratar alguém como sujeito de direitos implica reconhecer sua capacidade de compreender e questionar as decisões que lhe dizem respeito.

\subsection{O Desafio dos Modelos de Caixa-Preta}
\label{subsec:caixa-preta}

Um obstáculo técnico significativo ao direito à explicação é a natureza de ``caixa-preta'' de muitos sistemas de IA, especialmente os baseados em redes neurais profundas. Esses modelos atingem alto desempenho preditivo precisamente porque sua complexidade interna é tão grande que nem mesmo seus criadores conseguem explicar, em linguagem compreensível por seres humanos, como uma determinada decisão foi alcançada.

Esse fenômeno, conhecido como o trade-off entre acurácia e explicabilidade, coloca um dilema para o direito: se a proteção efetiva dos direitos fundamentais exige decisões explicáveis, e se os sistemas mais precisos são justamente os menos explicáveis, como conciliar esses dois imperativos?

A literatura jurídica tem proposto diferentes soluções para esse dilema. Uma primeira abordagem defende a proibição do uso de modelos de caixa-preta em decisões de alto risco, obrigando que sistemas utilizados em contextos penal, sanitário, creditício e laboratorial sejam necessariamente interpretáveis. Uma segunda abordagem aposta no desenvolvimento de técnicas de IA explicável (XAI — \textit{eXplainable Artificial Intelligence}) que permitam gerar explicações \textit{post hoc} mesmo para modelos complexos. Uma terceira abordagem, mais moderada, admite a utilização de modelos complexos desde que acompanhados de salvaguardas processuais robustas, como auditoria independente, testes de viés e supervisão humana efetiva \cite{Goodman2017EU, Selbst2018TIR, DoshiVelez2017MAR}.

A \MH{}, sem entrar no mérito técnico do debate, oferece critérios éticos para sua resolução: ``quando houver conflito entre a eficiência de um sistema e a proteção da dignidade humana, esta deve sempre prevalecer'' \cite{LeaoXIV2026MH}. Esse critério, se acolhido pelo ordenamento jurídico, implica que a explicabilidade deve ser considerada um requisito mínimo para a legitimidade do uso de IA em decisões que afetem direitos fundamentais.

\section{Discriminação Algorítmica}
\label{sec:discriminacao}

\subsection{Os Mecanismos da Discriminação Algorítmica}
\label{subsec:mecanismos-discriminacao}

A discriminação algorítmica ocorre quando sistemas de IA produzem resultados que desfavorecem sistematicamente determinados grupos ou indivíduos com base em características protegidas como raça, gênero, idade, orientação sexual, religião ou origem geográfica. Esse fenômeno pode decorrer de múltiplos fatores, que a literatura tem classificado em três categorias principais \cite{Noble2018AFS, Barocas2016BFF}.

A \textbf{discriminação por dados de treinamento} ocorre quando o conjunto de dados utilizado para treinar o modelo reflete preconceitos históricos existentes na sociedade. Se um algoritmo de recrutamento for treinado com dados históricos de contratação de uma empresa que discriminou mulheres, ele tenderá a reproduzir esse viés, excluindo candidatas do sexo feminino independentemente de sua qualificação.

A \textbf{discriminação por codificação} ocorre quando os próprios critérios e variáveis escolhidos para o modelo refletem escolhas valorativas que prejudicam determinados grupos. Por exemplo, um algoritmo de predição de reincidência criminal que utilize o histórico de prisões como variável preditiva tende a desfavorecer comunidades negras, que historicamente são alvo de policiamento mais intensivo e seletivo.

A \textbf{discriminação por \textit{proxy}} ocorre quando o algoritmo utiliza variáveis aparentemente neutras que funcionam como substitutas (\textit{proxies}) de características protegidas. O CEP de residência, por exemplo, pode funcionar como \textit{proxy} de raça ou classe social em muitas cidades brasileiras, produzindo discriminação indireta mesmo quando o modelo não utiliza explicitamente essas variáveis proibidas.

\subsection{A Discriminação Algorítmica na \MH{}}
\label{subsec:discriminacao-mh}

A \MH{} aborda a discriminação algorítmica como uma violação do princípio da igualdade e da justiça social. A Encíclica denuncia que a IA ``pode perpetuar e até ampliar as desigualdades existentes, criando um círculo vicioso no qual os já marginalizados são ainda mais prejudicados por sistemas que se apresentam como neutros e objetivos'' \cite{LeaoXIV2026MH}.

O documento estabelece uma conexão direta entre discriminação algorítmica e exclusão digital. Aqueles que não têm acesso à tecnologia, ou que não compreendem seu funcionamento, são duplamente vulneráveis: ficam excluídos dos benefícios da IA ao mesmo tempo em que podem ser prejudicados por suas decisões sem ter consciência disso ou meios de defesa. Essa constatação fundamenta a exigência de que a regulação da IA seja acompanhada de políticas de inclusão digital e de letramento algorítmico.

\subsection{Enfrentamento Jurídico da Discriminação Algorítmica no Brasil}
\label{subsec:enfrentamento-juridico-brasil}

O ordenamento jurídico brasileiro dispõe de instrumentos para o enfrentamento da discriminação algorítmica, embora ainda careça de um marco regulatório específico para a IA. O princípio da igualdade, inscrito no caput do art. 5º da Constituição Federal, combinado com a proibição de discriminação em razão de origem, raça, sexo, cor, idade e quaisquer outras formas de discriminação, constitui o fundamento primeiro para a tutela contra a discriminação algorítmica \cite{Brasil1988CF}.

A Lei de Igualdade Racial (Lei nº 12.288/2010), o Estatuto da Pessoa com Deficiência (Lei nº 13.146/2015) e a Lei de Combate à Discriminação (Lei nº 7.716/1989) oferecem instrumentos específicos de proteção. A LGPD, por sua vez, proíbe o tratamento de dados pessoais sensíveis para fins discriminatórios (art. 11) e determina que a revisão de decisões automatizadas deve levar em conta a não discriminação (art. 20, §2º) \cite{Brasil2018LGPD}.

No entanto, a efetividade desses instrumentos depende de adaptações interpretativas para o contexto algorítmico. A discriminação algorítmica frequentemente opera de forma indireta e difusa, tornando difícil a demonstração do nexo causal entre o funcionamento do sistema e o resultado discriminatório. A doutrina tem proposto, nesse sentido, a inversão do ônus da prova em casos de suspeita de discriminação algorítmica, bem como a instituição de auditorias obrigatórias de viés para sistemas de IA de alto risco \cite{Wachter2017TSP, Selbst2018TIR}.

\section{Vigilância, Controle e Privacidade}
\label{sec:vigilancia}

\subsection{O Capitalismo de Vigilância}
\label{subsec:capitalismo-vigilancia}

O conceito de capitalismo de vigilância, formulado por Shoshana Zuboff, descreve o modelo econômico emergente no qual a experiência humana é convertida em matéria-prima para a predição e o controle comportamental. Nesse modelo, os dados pessoais não são meramente um subproduto da atividade digital, mas o próprio objeto do processo produtivo: empresas de tecnologia extraem, analisam e monetizam dados comportamentais em escala massiva, utilizando algoritmos de IA para predizer e influenciar o comportamento futuro dos usuários \cite{Zuboff2019SC}.

A \MH{} dialoga com essa crítica ao denunciar que ``a coleta massiva de dados pessoais e sua utilização por sistemas algorítmicos cria uma assimetria de poder sem precedentes, na qual as pessoas são observadas, classificadas e manipuladas sem que tenham pleno conhecimento ou controle sobre esse processo'' \cite{LeaoXIV2026MH}.

O capitalismo de vigilância desafia a proteção jurídica da personalidade em pelo menos três dimensões. Em primeiro lugar, compromete a privacidade informacional, na medida em que dados pessoais são coletados e processados sem consentimento informado ou para finalidades diversas daquelas que motivaram a coleta. Em segundo lugar, afeta a autodeterminação informativa — o direito de controlar as próprias informações e decidir sobre sua circulação. Em terceiro lugar, ameaça a liberdade de escolha, pois a manipulação comportamental baseada em perfis algorítmicos pode induzir decisões que não refletem genuinamente a vontade do indivíduo.

\subsection{Vigilância Digital e Estado de Direito}
\label{subsec:vigilancia-estado-direito}

A utilização de sistemas de IA para fins de vigilância estatal suscita questões ainda mais graves do ponto de vista do Estado de Direito. Sistemas de reconhecimento facial em espaços públicos, análise preditiva de comportamento criminal, monitoramento massivo de comunicações e perfilamento ideológico são tecnologias que, se não submetidas a controles rigorosos, podem converter o Estado em Leviatã algorítmico.

O Supremo Tribunal Federal tem se manifestado crescentemente sobre essas questões. No julgamento da ADPF nº 695 (caso do reconhecimento facial no Metrô de São Paulo), o STF determinou que a utilização de tecnologias de reconhecimento facial pelo poder público exige lei específica, proporcionalidade, transparência e mecanismos efetivos de controle \cite{STF2022ADPF695}. Esse entendimento alinha-se à jurisprudência da Corte Europeia de Direitos Humanos, que no caso \textit{S. e Marper v. Reino Unido} (2008) já estabelecia limites rigorosos à retenção de dados biométricos pelo Estado.

A \MH{} contribui com essa discussão ao afirmar que ``a segurança pública não pode ser invocada como pretexto para a eliminação da privacidade e da liberdade'' \cite{LeaoXIV2026MH}. A Encíclica rejeita o falso dilema entre segurança e liberdade, sustentando que uma sociedade que sacrifica a privacidade em nome da segurança acaba por comprometer ambas.

\section{Identidade Pessoal e Subjetividade}
\label{sec:identidade}

\subsection{A Fragmentação da Identidade na Era Digital}
\label{subsec:fragmentacao-identidade}

A IA desafia também a integridade da identidade pessoal. Diferentemente das sociedades pré-digitais, onde a identidade de uma pessoa era construída principalmente por meio de narrativas biográficas e relações intersubjetivas, na era digital a identidade é cada vez mais definida por perfis algorítmicos construídos a partir de dados fragmentários — históricos de navegação, transações, localizações geográficas, interações em redes sociais.

Esses perfis não são neutros: eles refletem os interesses e as visões de mundo das empresas que os constroem. O perfil algorítmico de uma pessoa pode ser radicalmente diferente de sua autoimagem, e no entanto é esse perfil que determina seu acesso a crédito, emprego, seguros e serviços públicos. Cria-se, assim, uma tensão entre a identidade narrada (a história que a pessoa conta sobre si mesma) e a identidade algorítmica (o perfil que os sistemas atribuem à pessoa).

A \MH{} reconhece essa tensão ao afirmar que ``a pessoa humana não pode ser reduzida a um perfil ou a um conjunto de dados; ela é sempre mais do que aquilo que as máquinas podem capturar sobre ela'' \cite{LeaoXIV2026MH}. Esse reconhecimento fundamenta, no plano jurídico, o direito de não ser submetido a decisões baseadas exclusivamente em perfis algorítmicos, bem como o direito de contestar e corrigir as informações que compõem o perfil.

\subsection{A Voz Sintética e a Aparência Digital}
\label{subsec:voz-sintetica}

Tecnologias de IA generativa tornaram possível a criação de vozes sintéticas, imagens e vídeos hiper-realistas (deepfakes) que imitam pessoas reais com precisão perturbadora. Essas tecnologias abrem possibilidades criativas, mas também criam riscos graves para a identidade pessoal.

A clonagem não autorizada da voz ou da imagem de uma pessoa viola seus direitos de personalidade, especialmente o direito à própria imagem (art. 5º, V e X, da CF/88; arts. 12 a 20 do Código Civil). A \MH{} adverte que ``a criação de representações artificiais de pessoas reais, sem seu consentimento e especialmente com o intuito de enganar ou manipular, constitui uma forma de violência simbólica contra a pessoa, que perde o controle sobre sua própria representação'' \cite{LeaoXIV2026MH}.

Além disso, os deepfakes representam uma ameaça à confiança social necessária para o funcionamento das instituições democráticas. Vídeos fraudulentos de autoridades públicas podem ser utilizados para disseminar desinformação, influenciar eleições e desestabilizar governos. O direito à identidade pessoal, nesse contexto, adquire uma dimensão coletiva: proteger a integridade da representação de cada pessoa é também proteger a possibilidade de uma comunicação social confiável \cite{Chesney2019DF, Pariser2022THP}.

\section{Transumanismo, Pós-Humanismo e os Limites do Humano}
\label{sec:agencia-moral}

\subsection{O Projeto Transumanista}
\label{subsec:projeto-transumanista}

O transumanismo é um movimento intelectual e tecnológico que defende o uso da tecnologia para aprimorar as capacidades humanas para além de seus limites biológicos naturais, visando em última instância a superação da própria condição humana. Seus proponentes advogam a integração cada vez mais estreita entre seres humanos e máquinas, a extensão radical da longevidade, o aumento da inteligência por meios artificiais e, em suas versões mais radicais, o upload da consciência humana para suportes digitais \cite{Bostrom2005TH, Kurzweil2005THS}.

A \MH{} dedica atenção considerável ao transumanismo, identificando-o como uma das expressões mais extremas do paradigma tecnocrático. A Encíclica adverte que ``o desejo de superar os limites da condição humana, se não for acompanhado por uma reflexão profunda sobre o sentido da existência e sobre o bem integral da pessoa, pode levar a uma forma de desumanização mais sutil e perigosa do que as abertamente desumanizadoras'' \cite{LeaoXIV2026MH}.

A crítica da \MH{} ao transumanismo não se confunde com uma rejeição do progresso tecnológico. A Encíclica reconhece que intervenções tecnológicas no corpo e na mente humanos podem ser legítimas quando orientadas para a cura de patologias ou para a melhoria da qualidade de vida. A objeção recai sobre a pretensão de transcender a natureza humana como tal, substituindo a pessoa por uma entidade pós-humana que teria perdido precisamente aquilo que a torna humana: a corporeidade, a finitude, a historicidade e a capacidade de sofrer.

\subsection{Implicações Jurídicas do Transumanismo}
\label{subsec:implicacoes-juridicas}

O transumanismo levanta questões jurídicas profundas que desafiam categorias fundamentais do direito. Se a pessoa humana puder ser radicalmente modificada por intervenções tecnológicas, o que acontece com a noção de dignidade humana? Se a consciência puder ser transferida para suportes digitais, quem é o titular dos direitos de personalidade? Se a longevidade for estendida por séculos, como se reorganizam as relações contratuais, familiares e sucessórias?

A \MH{} não oferece respostas prontas a essas questões, mas estabelece um princípio orientador: ``a tecnologia deve servir à pessoa humana, e não o contrário; qualquer desenvolvimento tecnológico que reduza ou comprometa a humanidade do ser humano deve ser rejeitado'' \cite{LeaoXIV2026MH}. Esse princípio, se incorporado ao discurso jurídico, fundamenta a proibição de determinadas intervenções tecnológicas no corpo e na mente humanos, bem como a exigência de que o desenvolvimento tecnológico seja orientado por uma avaliação cuidadosa de seus impactos antropológicos.

\section{Síntese do Capítulo}
\label{sec:sumario-cap3}

Os desafios examinados neste capítulo demonstram que a IA não é neutra do ponto de vista antropológico e jurídico. Cada inovação tecnológica carrega consigo pressupostos sobre o que é a pessoa humana, sobre o que é o conhecimento, sobre o que é a decisão correta. Ignorar esses pressupostos é entregar à técnica a prerrogativa de definir, na prática, o conteúdo dos direitos fundamentais.

A \MH{}, ao explicitar os fundamentos antropológicos e éticos que devem orientar o desenvolvimento da IA, oferece critérios para que o direito possa exercer sua função de proteção da pessoa sem recorrer a uma rejeição simplista da tecnologia. Cinco desafios principais foram identificados:

\begin{enumerate}
    \item A \textbf{automação decisória} e seus impactos sobre a autonomia humana, que exige a afirmação do direito à explicação e à revisão humana das decisões automatizadas.
    \item A \textbf{discriminação algorítmica}, que reproduz e amplia desigualdades existentes e demanda mecanismos robustos de auditoria e não discriminação.
    \item A \textbf{vigilância e a perda de privacidade}, que criam assimetrias de poder incompatíveis com o Estado Democrático de Direito.
    \item A \textbf{fragmentação da identidade pessoal}, que impõe a necessidade de proteger a integridade da representação algorítmica da pessoa.
    \item O \textbf{transumanismo} e seus desafios aos limites do humano, que exigem uma reflexão sobre os limites éticos e jurídicos das intervenções tecnológicas na pessoa.
\end{enumerate}

O capítulo seguinte examina como o ordenamento jurídico brasileiro tem enfrentado esses desafios, identificando as potencialidades e as insuficiências dos instrumentos normativos existentes para a proteção da personalidade na era da IA.

% ============================================================
% CAPÍTULO 4: PROTEÇÃO JURÍDICA DA PERSONALIDADE 
%              NO DIREITO BRASILEIRO
% ============================================================

\chapter{Proteção Jurídica da Personalidade no Direito Brasileiro}
\label{cap:protecao-juridica}

\section{Introdução ao Capítulo}
\label{sec:intro-cap4}

O ordenamento jurídico brasileiro dispõe de um conjunto significativo de instrumentos normativos para a proteção da personalidade humana. A Constituição Federal de 1988, o Código Civil, a Lei Geral de Proteção de Dados Pessoais, o Marco Civil da Internet e a jurisprudência dos tribunais superiores compõem um arcabouço que, embora não tenha sido concebido especificamente para enfrentar os desafios da inteligência artificial, oferece categorias e princípios aplicáveis a esse novo contexto.

Este capítulo examina esses instrumentos, avaliando suas potencialidades e insuficiências para a proteção da personalidade na era da IA. O percurso é o seguinte: parte-se dos direitos da personalidade no Código Civil (\ref{sec:direitos-personalidade-cc}), analisa-se o tratamento constitucional da dignidade humana (\ref{sec:dignidade-constitucional}), examina-se a LGPD como marco regulatório da autodeterminação informativa (\ref{sec:lgpd}), avalia-se o Marco Civil da Internet (\ref{sec:mci}), e finalmente examina-se o Projeto de Lei nº 2.338/2023 como proposta de regulação específica da IA no Brasil (\ref{sec:pl-2338}).

\section{Os Direitos da Personalidade no Código Civil}
\label{sec:direitos-personalidade-cc}

\subsection{Fundamentos e Características}
\label{subsec:fundamentos-direitos-personalidade}

O Código Civil de 2002 dedicou, pela primeira vez na história legislativa brasileira, um capítulo autônomo aos direitos da personalidade (arts. 11 a 21). Essa inovação representou um avanço significativo na proteção jurídica da pessoa humana, superando a concepção meramente patrimonialista que predominava no Código Civil de 1916.

Os direitos da personalidade são definidos pela doutrina como aqueles direitos subjetivos que têm por objeto os atributos físicos, psíquicos e morais da pessoa, tutelando sua integridade física, intelectual e moral \cite{Amaral2003DBC}. Suas características fundamentais são a intransmissibilidade, a irrenunciabilidade, a indisponibilidade e a imprescritibilidade. O art. 11 do CC estabelece que ``com exceção dos casos previstos em lei, os direitos da personalidade são intransmissíveis e irrenunciáveis, não podendo o seu exercício sofrer limitação voluntária'' \cite{Brasil2002CC}.

O rol de direitos da personalidade no CC é exemplificativo, não taxativo. Estão expressamente previstos: a proteção ao nome (art. 16 a 19), à imagem (art. 20), à privacidade (art. 21), à integridade física e psíquica, e à honra. A doutrina reconhece, contudo, a existência de direitos da personalidade não expressamente codificados, como o direito à identidade pessoal, à autodeterminação informativa e ao esquecimento (este último objeto de intenso debate jurisprudencial).

\subsection{A Insuficiência do Modelo Civilista para os Desafios da IA}
\label{subsec:insuficiencia-civilista}

Embora os direitos da personalidade previstos no CC constituam um ponto de partida importante, sua concepção clássica mostra-se insuficiente para enfrentar os desafios específicos postos pela IA. Essa insuficiência manifesta-se em pelo menos três planos.

Em \textbf{primeiro lugar}, o modelo civilista tradicional é essencialmente reativo e reparatório: pressupõe a ocorrência de um dano concreto para que o titular possa exercer seu direito de exigir reparação. A IA, contudo, produz frequentemente danos difusos e sistêmicos que não se enquadram facilmente no esquema tradicional da responsabilidade civil (dano, nexo causal, culpa). A discriminação algorítmica, por exemplo, afeta grupos inteiros de forma indireta, tornando difícil a individualização do dano e a identificação do nexo causal.

Em \textbf{segundo lugar}, o direito da personalidade clássico foi concebido para tutelar a pessoa contra agressões identificáveis e localizáveis. A opacidade dos sistemas algorítmicos torna difícil sequer saber se e quando uma violação ocorreu. Como exercer o direito de exigir reparação se a própria existência do dano é encoberta pela caixa-preta algorítmica?

Em \textbf{terceiro lugar}, os direitos da personalidade foram pensados primordialmente como direitos de defesa contra o Estado (dimensão vertical), mas os principais riscos decorrentes da IA emanam de atores privados, grandes corporações de tecnologia, que operam em escala global e com poder econômico frequentemente superior ao de muitos Estados nacionais. Essa dimensão horizontal dos direitos fundamentais (eficácia privada ou horizontal) exige uma reconstrução teórica dos instrumentos de proteção.

\section{A Dignidade da Pessoa Humana como Fundamento Constitucional}
\label{sec:dignidade-constitucional}

\subsection{A Dignidade como Princípio Estruturante}
\label{subsec:dignidade-principio}

A Constituição Federal de 1988 elegeu a dignidade da pessoa humana como fundamento da República (art. 1º, III) e como princípio informador de todo o ordenamento jurídico. A doutrina constitucional brasileira, especialmente a partir das contribuições de Ingo Sarlet e Daniel Sarmento, tem desenvolvido um conceito de dignidade que vai além de sua formulação meramente negativa (proibição de tratamento degradante) para abranger também uma dimensão positiva de promoção das condições para o pleno desenvolvimento da personalidade \cite{Sarlet2021DF, Sarmento2016DPF}.

A dignidade da pessoa humana, nessa perspectiva, não é um direito entre outros, mas o fundamento axiológico que dá unidade e sentido ao sistema de direitos fundamentais. Ela opera como limite (proibindo intervenções degradantes), como dever de proteção (obrigando o Estado a proteger a pessoa contra agressões de terceiros) e como mandato de promoção (exigindo políticas públicas que criem as condições materiais para o desenvolvimento da personalidade).

Essas três dimensões são relevantes para a regulação da IA. Como \textbf{limite}, a dignidade proíbe a utilização de pessoas como meros meios para fins econômicos ou tecnológicos, vedando práticas como a manipulação algorítmica não consentida ou a redução da pessoa a um perfil de dados. Como \textbf{dever de proteção}, ela obriga o Estado a editar normas que protejam as pessoas contra os riscos da IA e a fiscalizar sua observância. Como \textbf{mandato de promoção}, ela exige políticas de inclusão digital, letramento algorítmico e acesso equitativo às tecnologias.

\subsection{A Dignidade e a Proteção Contra Riscos Algorítmicos}
\label{subsec:dignidade-riscos}

A aplicação do princípio da dignidade à proteção contra riscos algorítmicos encontra amparo na jurisprudência do Supremo Tribunal Federal. No julgamento da ADI nº 6.640 (que tratou da validade constitucional da LGPD), o STF reconheceu que a proteção de dados pessoais é expressão do direito fundamental à privacidade e à autodeterminação informativa, ambos ancorados no princípio da dignidade da pessoa humana \cite{STF2020ADI6640}.

A \MH{} oferece uma contribuição relevante a esse debate ao articular o princípio da dignidade com a crítica ao paradigma tecnocrático. Para a Encíclica, a dignidade não é apenas um conceito jurídico abstrato, mas uma realidade concreta que pode ser violada sempre que a pessoa for tratada como objeto, como meio ou como dado. A proteção contra a redução algorítmica da pessoa a um perfil ou a uma classificação é, portanto, uma exigência da própria dignidade \cite{LeaoXIV2026MH}.

\section{A Lei Geral de Proteção de Dados Pessoais}
\label{sec:lgpd}

\subsection{Estrutura e Princípios da LGPD}
\label{subsec:estrutura-lgpd}

A Lei nº 13.709/2018, Lei Geral de Proteção de Dados Pessoais, representa o marco normativo mais relevante para a proteção da personalidade na era digital no Brasil. Inspirada no Regulamento Geral de Proteção de Dados da União Europeia (GDPR), a LGPD estabelece um conjunto abrangente de regras para o tratamento de dados pessoais por pessoas naturais e jurídicas, públicas e privadas.

A LGPD se estrutura em torno de princípios fundamentais que orientam toda a atividade de tratamento de dados: finalidade, adequação, necessidade, livre acesso, qualidade dos dados, transparência, segurança, prevenção, não discriminação e responsabilização (art. 6º). Esses princípios funcionam como cláusulas gerais que devem informar a interpretação e aplicação de todas as disposições da Lei.

Para os fins deste trabalho, destacam-se três princípios especialmente relevantes:

O \textbf{princípio da necessidade} (art. 6º, III) determina que o tratamento de dados deve limitar-se ao mínimo necessário para a realização de suas finalidades. Esse princípio é particularmente importante para conter a coleta massiva de dados que alimenta os sistemas de IA, estabelecendo um limite à voracidade informacional do capitalismo de vigilância.

O \textbf{princípio da não discriminação} (art. 6º, IX) proíbe a utilização de dados pessoais para fins discriminatórios ilícitos ou abusivos. Esse princípio constitui um fundamento normativo direto para o combate à discriminação algorítmica.

O \textbf{princípio da responsabilização e prestação de contas} (art. 6º, X) exige que o controlador demonstre a adoção de medidas eficazes capazes de comprovar a observância das normas de proteção de dados. Esse princípio fundamenta a exigência de auditorias e avaliações de impacto à proteção de dados, essenciais para o controle preventivo de riscos algorítmicos.

\subsection{O Direito à Revisão de Decisões Automatizadas}
\label{subsec:revisao-automatizada}

O art. 20 da LGPD consagra o direito do titular de dados pessoais ``a solicitar revisão de decisões tomadas exclusivamente com base em tratamento automatizado de dados pessoais que afetem seus interesses, incluídas as decisões destinadas a definir o seu perfil pessoal, profissional, de consumo e de crédito ou os aspectos de sua personalidade''.

Esse dispositivo é inovador no direito brasileiro, pois reconhece que decisões algorítmicas podem afetar a personalidade humana e cria um mecanismo processual para sua contestação. O §1º do mesmo artigo determina que a revisão deve ser realizada por pessoa natural, estabelecendo a garantia de supervisão humana como contrapeso à automação decisória.

Contudo, a implementação prática desse direito enfrenta desafios significativos. Estudos apontam que o exercício do direito de revisão é frequentemente dificultado pela falta de transparência dos sistemas algorítmicos, pela assimetria de informações entre controlador e titular, e pela ausência de padrões claros para a realização da revisão \cite{Bioni2021PPD, Doneda2022DPD}. Além disso, a redação do art. 20 refere-se a decisões ``tomadas exclusivamente com base em tratamento automatizado'', o que exige um elevado standard probatório por parte do titular que deseja contestar a decisão.

\subsection{As Avaliações de Impacto Algorítmico}
\label{subsec:avaliacoes-impacto}

Um dos instrumentos mais promissores para a proteção preventiva contra riscos algorítmicos é a Avaliação de Impacto à Proteção de Dados Pessoais (AIPD), prevista no art. 38 da LGPD. A AIPD é um processo sistemático de identificação, análise e mitigação de riscos à privacidade e à proteção de dados associados a operações de tratamento.

A Autoridade Nacional de Proteção de Dados (ANPD) tem incentivado a realização de AIPDs especialmente para sistemas de IA de alto risco. A Resolução ANPD nº 4/2024 estabeleceu diretrizes para a elaboração de relatórios de impacto, incluindo a descrição do funcionamento do sistema, a identificação de riscos e as medidas de mitigação adotadas.

A doutrina tem proposto a ampliação desse instrumento para além da mera proteção de dados, transformando-o em uma verdadeira Avaliação de Impacto Algorítmico (AIA) que considere não apenas riscos à privacidade, mas também riscos de discriminação, exclusão, manipulação e violação de direitos fundamentais \cite{Doneda2022DPD, Edwards2017PC}. A \MH{} oferece fundamentos éticos para essa ampliação, ao insistir que a proteção da pessoa humana deve ser abrangente e não restrita a uma única dimensão de sua personalidade \cite{LeaoXIV2026MH}.

\section{O Marco Civil da Internet}
\label{sec:mci}

\subsection{Princípios e Direitos}
\label{subsec:principios-mci}

A Lei nº 12.965/2014, Marco Civil da Internet (MCI), estabelece princípios, garantias, direitos e deveres para o uso da Internet no Brasil. Embora não trate especificamente de IA, o MCI oferece fundamentos relevantes para a proteção da personalidade no ambiente digital.

O MCI consagra, em seu art. 3º, os princípios de neutralidade da rede, privacidade, proteção de dados pessoais, liberdade de expressão e responsabilização dos agentes. Esses princípios informam a interpretação de todo o ordenamento digital e devem orientar também a regulação da IA.

Especialmente relevante é o art. 10 do MCI, que determina a proteção dos registros de conexão e de acesso a aplicações de internet, ressalvadas as hipóteses de guarda obrigatória e de requisição judicial. Esse dispositivo estabelece um regime de proteção de dados que se aplica também aos dados utilizados por sistemas de IA.

\subsection{A Responsabilidade Civil dos Provedores}
\label{subsec:responsabilidade-mci}

O MCI estabelece um regime especial de responsabilidade civil para provedores de aplicação. O art. 19 determina que o provedor só pode ser responsabilizado por danos decorrentes de conteúdo gerado por terceiros se, após ordem judicial específica, não tomar as providências para tornar indisponível o conteúdo. Esse regime tem sido objeto de intenso debate no Supremo Tribunal Federal (RE nº 1.037.396/SP e RE nº 1.057.258/SP), que discute sua constitucionalidade e extensão.

Para os sistemas de IA gerativa, o regime do art. 19 mostra-se particularmente problemático. Se um modelo de IA produz conteúdo danoso (deepfake difamatório, informação falsa, discurso de ódio), aplica-se a responsabilidade do provedor? A resposta não é clara, pois o conteúdo não é propriamente ``gerado por terceiros'' (o usuário que solicitou a geração) nem pelo provedor (que apenas disponibiliza o modelo). Essa lacuna normativa evidencia a necessidade de um marco regulatório específico para a IA.

\section{O Projeto de Lei nº 2.338/2023}
\label{sec:pl-2338}

\subsection{Contexto e Conteúdo da Proposta}
\label{subsec:contexto-pl}

O Projeto de Lei nº 2.338/2023, em tramitação no Congresso Nacional, propõe um marco regulatório abrangente para a inteligência artificial no Brasil. Inspirado no \textit{AI Act} da União Europeia, o PL adota uma abordagem baseada em risco, classificando os sistemas de IA em diferentes níveis de risco (risco excessivo, alto risco, risco limitado e risco mínimo) e estabelecendo obrigações proporcionais ao nível de risco.

Para os fins deste trabalho, destacam-se três aspectos do PL.

Em primeiro lugar, o PL proíbe determinadas práticas consideradas de risco excessivo, como sistemas de pontuação social, reconhecimento emocional em ambientes de trabalho e educacionais, e categorização biométrica baseada em características sensíveis. Essas proibições alinham-se com as preocupações da \MH{} sobre a instrumentalização da pessoa humana pela tecnologia.

Em segundo lugar, o PL estabelece um conjunto de obrigações para sistemas de IA de alto risco, incluindo avaliação de impacto algorítmico, transparência, supervisão humana, documentação técnica e governança de dados. Essas obrigações concretizam, no plano normativo, muitos dos princípios éticos propostos pela \MH{}, especialmente transparência e responsabilidade.

Em terceiro lugar, o PL cria o Sistema Nacional de Regulação de Inteligência Artificial (SINRIA) como estrutura de governança com competência para expedir normas, fiscalizar e sancionar infrações. A eficácia do PL dependerá crucialmente da capacidade do SINRIA de exercer essas funções de forma independente, técnica e eficiente.

\subsection{Diálogo com a \MH{}}
\label{subsec:dialogo-pl-mh}

O PL 2.338/2023 e a \MH{} operam em registros distintos (o primeiro no âmbito da regulação técnica e jurídica, a segunda no campo da reflexão ética e antropológica), mas ambos convergem na preocupação fundamental com a proteção da pessoa humana diante dos riscos da IA.

A \MH{} oferece ao debate regulatório brasileiro uma fundamentação antropológica que o PL, por sua própria natureza técnica, não desenvolve. Enquanto o PL estabelece o que deve ser regulado (classificação de riscos) e como (obrigações, sanções, governança), a Encíclica explica por que a proteção da pessoa é necessária e quais valores estão em jogo.

Essa complementaridade é particularmente evidente em três pontos. Primeiro, a classificação de riscos adotada pelo PL encontra correspondência na gradação ética proposta pela \MH{}, que distingue entre usos aceitáveis, problemáticos e inaceitáveis da IA. Segundo, as obrigações de transparência e supervisão humana do PL ecoam a exigência ética de que a pessoa não seja submetida a decisões que não possa compreender ou controlar. Terceiro, a estrutura de governança do PL pode ser vista como uma concretização institucional do princípio de subsidiariedade, na medida em que estabelece uma instância pública de coordenação e fiscalização sem eliminar a responsabilidade dos atores privados.

\section{Síntese do Capítulo}
\label{sec:sumario-cap4}

O exame dos instrumentos jurídicos brasileiros para a proteção da personalidade revela um panorama misto. Por um lado, o ordenamento dispõe de fundamentos constitucionais sólidos (dignidade da pessoa humana), institutos civilistas consolidados (direitos da personalidade), um marco regulatório de proteção de dados em desenvolvimento (LGPD) e uma proposta de regulação específica da IA (PL 2.338/2023). Por outro lado, identificam-se insuficiências significativas:

\begin{enumerate}
    \item O \textbf{modelo civilista clássico}, de caráter reativo e reparatório, mostra-se inadequado para lidar com danos difusos e sistêmicos produzidos por sistemas algorítmicos.
    \item A \textbf{LGPD}, apesar de seu avanço, carece de instrumentos específicos para o controle da discriminação algorítmica e para a garantia efetiva do direito à explicação.
    \item O \textbf{MCI} não oferece resposta adequada para a responsabilidade civil por conteúdos gerados por IA.
    \item O \textbf{PL 2.338/2023}, embora promissor, ainda está em tramitação e seu texto final é incerto.
\end{enumerate}

Diante desse quadro, o capítulo seguinte propõe-se a extrair das contribuições da \MH{} orientações para o aperfeiçoamento do regime jurídico de proteção da personalidade, especialmente no que diz respeito à introdução de princípios antropológicos explícitos na regulação da IA, ao fortalecimento dos mecanismos preventivos e à articulação entre os diferentes instrumentos normativos existentes.

% ============================================================
% CAPÍTULO 5 — CONTRIBUIÇÕES DA ENCÍCLICA MAGNIFICA HUMANITAS
%              PARA A REGULAÇÃO DA INTELIGÊNCIA ARTIFICIAL
% ============================================================

\chapter{Contribuições da Magnifica Humanitas para a Regulação da Inteligência Artificial}
\label{cap:contribuicoes}

\section{Introdução ao Capítulo}
\label{sec:intro-cap5}

Os capítulos precedentes estabeleceram as bases deste trabalho: a fundamentação antropológica e ética da \MH{} (capítulo 2), os desafios específicos que a IA impõe à personalidade humana (capítulo 3) e os instrumentos jurídicos brasileiros disponíveis para sua proteção (capítulo 4). Este quinto capítulo articula esses três eixos, propondo uma síntese das contribuições que a \MH{} pode oferecer para o aperfeiçoamento da regulação da inteligência artificial.

A tese central deste capítulo é que a \MH{}, embora não seja um documento jurídico nem proponha regras específicas, oferece uma fundamentação antropológica robusta que pode preencher lacunas significativas do debate regulatório contemporâneo. Essas contribuições organizam-se em três planos: a crítica ao paradigma tecnocrático (\ref{sec:critica-tecnocratico}), a formulação de princípios éticos para a governança algorítmica (\ref{sec:principios-eticos-mh}), e a articulação concreta com os instrumentos jurídicos brasileiros (\ref{sec:articulacao-juridica}).

\section{A Crítica ao Paradigma Tecnocrático}
\label{sec:critica-tecnocratico}

\subsection{O Paradigma Tecnocrático como Risco Sistêmico}
\label{subsec:paradigma-risco-sistemico}

Uma das contribuições mais originais da \MH{} para o debate sobre a regulação da IA é a formulação do conceito de ``paradigma tecnocrático''. Se o capítulo segundo examinou esse conceito em sua dimensão antropológica, importa agora explorar suas implicações para o direito e para a regulação.

O paradigma tecnocrático, como desenvolvido pela \MH{} a partir de uma reflexão já presente na encíclica \textit{Laudato Si'} (§106–114), designa não apenas um conjunto de tecnologias, mas uma visão de mundo que reduz a realidade a seus aspectos mensuráveis e manipuláveis, subordinando todas as dimensões da existência humana à lógica da eficiência técnica e do lucro \cite{Francisco2015LS, LeaoXIV2026MH}.

Para a regulação da IA, a crítica ao paradigma tecnocrático tem consequências diretas. Em primeiro lugar, ela revela que os riscos da IA não são meramente acidentais ou decorrentes de má aplicação, mas podem ser sistêmicos, isto é, inerentes à própria lógica de desenvolvimento tecnológico quando dissociada de uma reflexão ética adequada. Em segundo lugar, ela adverte contra a tentação de resolver problemas éticos com soluções puramente técnicas (a chamada ``solução tecnocrática''), que frequentemente agravam os problemas que pretendem resolver.

Essa perspectiva oferece uma base crítica para avaliar propostas regulatórias que se limitam a estabelecer regras procedimentais sem questionar os pressupostos antropológicos do desenvolvimento tecnológico. Para a \MH{}, a regulação da IA não pode ser reduzida a uma questão de compliance técnica; ela deve envolver uma reflexão substantiva sobre que tipo de sociedade queremos construir e que lugar reservamos à pessoa humana.

\subsection{A Recusa da Neutralidade Tecnológica}
\label{subsec:recusa-neutralidade}

A \MH{} oferece também uma crítica à tese da neutralidade tecnológica — a ideia de que a tecnologia é intrinsecamente neutra e que seus efeitos dependem exclusivamente do uso que se faz dela. Para a Encíclica, essa tese ignora que as tecnologias incorporam valores, pressupostos e visões de mundo que moldam a experiência humana de forma não trivial \cite{LeaoXIV2026MH}.

No campo da regulação da IA, a recusa da neutralidade tecnológica tem implicações importantes. Se os sistemas de IA não são neutros — se eles carregam em seu próprio desenho técnico escolhas valorativas sobre o que é relevante, justo ou desejável —, então a regulação não pode limitar-se a controlar usos abusivos, mas deve também incidir sobre o próprio desenho dos sistemas.

\section{Princípios Éticos para a Governança Algorítmica}
\label{sec:principios-eticos-mh}

\subsection{Dignidade Humana como Critério Regulatório}
\label{subsec:dignidade-criterio}

O primeiro e mais fundamental princípio que a \MH{} oferece para a regulação da IA é a dignidade da pessoa humana como critério substantivo de avaliação e limite intransponível para o desenvolvimento tecnológico.

Aplicada à regulação algorítmica, a dignidade opera em três dimensões complementares: como limite negativo (proibindo usos da IA que reduzam a pessoa a objeto), como dever de proteção (exigindo que o Estado crie mecanismos eficazes de tutela) e como mandato de promoção (orientando políticas de desenvolvimento tecnológico que coloquem a pessoa no centro).

Na prática regulatória, essas três dimensões traduzem-se em exigências concretas: a proibição de sistemas de IA que instrumentalizem a pessoa sem seu consentimento informado; a obrigatoriedade de avaliações de impacto algorítmico prévias ao lançamento de sistemas de alto risco; e a previsão de mecanismos de supervisão humana efetiva que garantam o controle último da pessoa sobre as decisões que a afetam.

\subsection{Transparência e Direito à Explicação}
\label{subsec:transparencia-explicacao}

A \MH{} afirma, em consonância com o §125 que trata do direito à explicação, que ``a pessoa humana tem o direito de compreender os critérios que orientam as decisões que lhe dizem respeito'' \cite{LeaoXIV2026MH}. Esse princípio tem implicações diretas para a regulação da IA, especialmente no que tange às exigências de explicabilidade e transparência.

O direito à explicação, como desenvolvido pela \MH{}, não se confunde com a mera disponibilização de informações técnicas sobre o funcionamento do algoritmo. Ele exige que a explicação seja acessível, compreensível e significativa para a pessoa afetada pela decisão. Isso implica um dever positivo de traduzir o funcionamento técnico do sistema em linguagem acessível e de fornecer os fundamentos concretos da decisão particular.

Esse entendimento alinha-se com as propostas mais avançadas da doutrina especializada, que distinguem entre explicabilidade global (compreensão do funcionamento geral do sistema) e explicabilidade local (compreensão das razões de uma decisão particular) \cite{DoshiVelez2017TMI, Selbst2018INC}. A \MH{} oferece uma fundamentação antropológica para essa dupla exigência: a pessoa precisa compreender tanto a lógica geral dos sistemas que a afetam quanto as razões específicas de cada decisão que recai sobre si.

\subsection{Subsidiariedade e Governança Participativa}
\label{subsec:subsidiariedade-governanca}

O princípio da subsidiariedade, central na doutrina social da Igreja, oferece um critério relevante para a estruturação da governança algorítmica. A \MH{} (§105) recorda que ``nenhuma instância superior deve assumir funções que podem ser adequadamente desempenhadas pelas instâncias inferiores'' \cite{LeaoXIV2026MH}.

Aplicado à regulação da IA, o princípio da subsidiariedade implica que a regulação não deve excluir a responsabilidade dos próprios desenvolvedores e usuários de sistemas de IA. O Estado não deve substituir-se à iniciativa privada nem à sociedade civil na definição de padrões éticos e técnicos, mas deve criar as condições para que esses atores exerçam sua responsabilidade de forma efetiva.

Essa perspectiva sugere um modelo de governança algorítmica que articule diferentes níveis: autorregulação regulada (com padrões éticos definidos pelos próprios desenvolvedores, sob supervisão pública), codesign normativo (com participação de múltiplos stakeholders na definição de regras), e fiscalização pública independente (com sanções proporcionais para infrações).

\subsection{Solidariedade e Inclusão Digital}
\label{subsec:solidariedade-inclusao}

O princípio da solidariedade, também central na doutrina social da Igreja, orienta a regulação da IA para a superação das assimetrias e desigualdades que a tecnologia pode aprofundar. A \MH{} enfatiza que o desenvolvimento tecnológico deve beneficiar toda a família humana e não apenas alguns privilegiados (§189), e que a exclusão digital é uma forma contemporânea de exclusão social que viola a dignidade da pessoa \cite{LeaoXIV2026MH}.

Para a regulação da IA, a solidariedade traduz-se em exigências de inclusão digital, letramento algorítmico e acesso equitativo às tecnologias. Essas dimensões são frequentemente negligenciadas nos debates regulatórios, que tendem a concentrar-se nos riscos sem considerar as oportunidades e a distribuição desigual de benefícios.

\section{Articulação com os Instrumentos Jurídicos Brasileiros}
\label{sec:articulacao-juridica}

\subsection{Adequação e Aperfeiçoamento da LGPD}
\label{subsec:lgpd-aperfeicoamento}

A LGPD, como examinado no capítulo anterior, já contém instrumentos relevantes para a proteção da personalidade contra riscos algorítmicos, especialmente o direito à revisão de decisões automatizadas (art. 20) e a avaliação de impacto à proteção de dados (art. 38). Contudo, a perspectiva oferecida pela \MH{} permite identificar lacunas e propor aperfeiçoamentos.

Em primeiro lugar, o direito à revisão de decisões automatizadas poderia ser fortalecido com a explicitação de que a revisão deve ser realizada por pessoa natural com competência técnica e autonomia para modificar a decisão — e não mera formalidade de chancela automática. A \MH{} fundamenta essa exigência na necessidade de que a decisão final sobre a pessoa seja efetivamente humana.

Em segundo lugar, a avaliação de impacto poderia ser ampliada para incluir não apenas riscos à privacidade e proteção de dados, mas também riscos à dignidade, à não discriminação, à autonomia e ao desenvolvimento da personalidade. Trata-se de transformar a AIPD em uma verdadeira Avaliação de Impacto Algorítmico (AIA), inspirada no princípio de que todos os aspectos da personalidade merecem proteção.

Em terceiro lugar, a LGPD poderia prever mecanismos mais efetivos de transparência algorítmica, incluindo a obrigatoriedade de auditorias periódicas independentes para sistemas de IA de alto risco e a publicação de relatórios de impacto algorítmico em formato acessível.

\subsection{Contribuições para o PL 2.338/2023}
\label{subsec:contribuicoes-pl}

O PL 2.338/2023, atualmente em tramitação, representa uma oportunidade ímpar para incorporar, no direito positivo brasileiro, as contribuições da reflexão ética e antropológica da \MH{}.

Em primeiro lugar, o PL poderia explicitar, em seus fundamentos, a dignidade da pessoa humana como critério nuclear de avaliação dos sistemas de IA. Embora a dignidade já seja fundamento constitucional, sua explicitação no marco regulatório da IA teria o efeito de orientar a interpretação de todas as demais disposições e de sinalizar o compromisso do ordenamento com a proteção substantiva da pessoa — e não apenas com requisitos formais de compliance.

Em segundo lugar, o PL poderia fortalecer as exigências de transparência e supervisão humana para sistemas de IA que afetam direitos fundamentais, inspirando-se na distinção entre os três níveis de transparência propostos pela \MH{}: acesso à informação, compreensão do funcionamento e capacidade de contestação.

Em terceiro lugar, o PL poderia prever mecanismos de governança participativa que envolvam não apenas especialistas e representantes do setor, mas também a sociedade civil e as comunidades potencialmente afetadas pelos sistemas de IA. A subsidiariedade, como princípio de justiça social, exige que aqueles que serão afetados pelas decisões algorítmicas tenham voz nos processos decisórios que as definem.

\section{Síntese do Capítulo}
\label{sec:sumario-cap5}

A \MH{} oferece contribuições significativas para a regulação da IA em três planos. No plano da \textbf{crítica}, fornece instrumentos conceituais para questionar o paradigma tecnocrático e a tese da neutralidade tecnológica. No plano dos \textbf{princípios}, oferece fundamentos antropológicos robustos para a dignidade, a transparência, a subsidiariedade e a solidariedade como critérios regulatórios. No plano da \textbf{articulação concreta}, aponta direções para o aperfeiçoamento da LGPD e para o enriquecimento do PL 2.338/2023.

Essas contribuições não substituem o trabalho técnico da regulação, mas conferem-lhe uma fundamentação ética e antropológica que, como se argumentou ao longo do trabalho, é indispensável para que a regulação da IA não seja reduzida a um exercício de mera gestão de riscos, mas se constitua como verdadeira proteção da pessoa humana no contexto da transformação tecnológica.

% ============================================================
% CAPÍTULO 6: CONCLUSÃO
% ============================================================

\chapter{Conclusão}
\label{cap:conclusao}

\section{Retomada do Problema e das Hipóteses}
\label{sec:retomada-hipoteses}

Esta monografia partiu do seguinte problema de pesquisa: \textit{de que forma a Encíclica \MH{} pode contribuir para a proteção jurídica da personalidade humana diante dos desafios impostos pelo desenvolvimento e difusão da inteligência artificial?}

Para responder a essa questão, foram formuladas quatro hipóteses de trabalho:

\begin{itemize}
    \item \textbf{H1}: A \MH{} oferece uma fundamentação antropológica robusta para a proteção da personalidade, ancorada no personalismo cristão e em uma concepção relacional e transcendental da pessoa.
    \item \textbf{H2}: A IA apresenta riscos específicos e qualitativamente novos à personalidade humana, que transcendem as categorias tradicionais de proteção jurídica.
    \item \textbf{H3}: O ordenamento jurídico brasileiro dispõe de instrumentos relevantes para a proteção da personalidade, mas estes apresentam lacunas e insuficiências diante dos desafios específicos da IA.
    \item \textbf{H4}: A \MH{} oferece contribuições normativas que podem orientar o aperfeiçoamento da regulação da IA no Brasil, especialmente nos planos principiológico, crítico e institucional.
\end{itemize}

O percurso investigativo realizado nos capítulos precedentes permite agora avaliar em que medida essas hipóteses foram confirmadas.

\section{Verificação das Hipóteses}
\label{sec:verificacao-hipoteses}

\subsection{H1: A Fundamentação Antropológica da \MH{}}
\label{subsec:verificacao-h1}

A análise desenvolvida no capítulo segundo confirmou integralmente a primeira hipótese. A \MH{} oferece, de fato, uma fundamentação antropológica robusta para a proteção da personalidade, articulada em três dimensões: ontológica (a pessoa como ser dotado de dignidade intrínseca), relacional (a pessoa como ser constituído nas e pelas relações com os outros) e transcendental (a pessoa como ser aberto à transcendência).

O conceito de ``irredutibilidade da pessoa aos seus dados'', que a Encíclica desenvolve em diversos parágrafos (§3, §101, §186, §232), mostrou-se particularmente relevante para o debate contemporâneo sobre proteção de dados e regulação da IA. Esse conceito oferece um fundamento antropológico para a tese de que a pessoa não pode ser reduzida a seus traços digitais, a seus perfis algorítmicos ou a suas classificações preditivas.

Os cinco princípios éticos identificados (dignidade, subsidiariedade, solidariedade, transparência e responsabilidade) constituem um quadro normativo coerente e aplicável à regulação da IA.

\subsection{H2: Os Riscos Específicos da IA}
\label{subsec:verificacao-h2}

A análise desenvolvida no capítulo terceiro confirmou a segunda hipótese, demonstrando que a IA apresenta riscos qualitativamente novos à personalidade humana. A identificação de quatro dimensões de risco (a automação decisória e o direito à explicação, a discriminação algorítmica, o capitalismo de vigilância e a identidade pessoal em questão, e os desafios do transumanismo) evidenciou que a proteção da personalidade na era da IA não pode ser adequadamente enfrentada com os instrumentos tradicionais do direito civil e constitucional.

Particularmente significativa foi a constatação de que a opacidade dos sistemas algorítmicos desafia a própria possibilidade de exercício de direitos, na medida em que a pessoa afetada por uma decisão automatizada frequentemente não dispõe dos meios para sequer identificar a existência da violação.

\subsection{H3: As Insuficiências do Ordenamento Brasileiro}
\label{subsec:verificacao-h3}

A análise do capítulo quarto confirmou a terceira hipótese, demonstrando que, embora o ordenamento jurídico brasileiro disponha de instrumentos relevantes, estes apresentam lacunas significativas. A LGPD, apesar de seu caráter inovador, carece de instrumentos específicos para o controle da discriminação algorítmica e para a garantia efetiva do direito à explicação. O Marco Civil da Internet não oferece resposta adequada para a responsabilidade civil por conteúdos gerados por IA. O PL 2.338/2023, embora promissor, ainda está em tramitação e seu texto final é incerto.

A principal lacuna identificada foi a ausência, nos instrumentos jurídicos brasileiros, de uma fundamentação antropológica explícita que oriente a interpretação e aplicação das normas de proteção da pessoa diante dos riscos algorítmicos.

\subsection{H4: As Contribuições Regulatórias da \MH{}}
\label{subsec:verificacao-h4}

O capítulo quinto confirmou a quarta hipótese, demonstrando que a \MH{} oferece contribuições normativas significativas para o aperfeiçoamento da regulação da IA. No plano crítico, fornece instrumentos conceituais para questionar o paradigma tecnocrático e a tese da neutralidade tecnológica. No plano principiológico, oferece fundamentos antropológicos robustos para dignidade, transparência, subsidiariedade e solidariedade como critérios regulatórios. No plano institucional, aponta direções para o aperfeiçoamento da LGPD e para o enriquecimento do PL 2.338/2023.

A tese central, de que a \MH{} pode fundamentar uma regulação da IA que vá além da mera gestão de riscos e se constitua como verdadeira proteção da pessoa humana, foi confirmada pela análise.

\section{Resposta ao Problema de Pesquisa}
\label{sec:resposta-problema}

À luz das evidências e análises apresentadas, é possível oferecer uma resposta ao problema de pesquisa proposto.

A \MH{} contribui para a proteção jurídica da personalidade humana diante dos desafios da IA em três planos fundamentais. Em \textbf{primeiro lugar}, oferece uma fundamentação antropológica para os direitos da personalidade que supera as concepções meramente formalistas ou procedimentais predominantes no debate regulatório contemporâneo. A Encíclica recorda que a proteção da pessoa não é uma questão técnica a ser resolvida por mecanismos de compliance, mas uma exigência ética fundamental que deve orientar todo o desenvolvimento tecnológico.

Em \textbf{segundo lugar}, a \MH{} formula princípios éticos (dignidade, transparência, subsidiariedade, solidariedade, responsabilidade) que podem operar como critérios substantivos para a avaliação e o aperfeiçoamento dos instrumentos jurídicos de proteção da personalidade. Esses princípios não substituem o trabalho técnico da regulação, mas conferem-lhe uma orientação axiológica clara e consistente.

Em \textbf{terceiro lugar}, a \MH{} oferece uma crítica ao paradigma tecnocrático que permite compreender os riscos da IA não como fenômenos isolados e corrigíveis, mas como manifestações de uma visão de mundo que tende a reduzir a pessoa a seus aspectos mensuráveis e funcionais. Essa crítica fundamenta a exigência de que a regulação da IA não se limite a controlar abusos, mas promova ativamente um desenvolvimento tecnológico centrado na pessoa.

\section{Limitações da Pesquisa}
\label{sec:limitacoes}

Este trabalho apresenta limitações que devem ser explicitadas para a adequada compreensão de seus resultados.

Em \textbf{primeiro lugar}, a pesquisa limitou-se ao exame do ordenamento jurídico brasileiro, não abrangendo o direito comparado nem o direito internacional, salvo referências pontuais ao GDPR europeu. Uma investigação mais ampla poderia identificar contribuições adicionais da \MH{} para outros sistemas jurídicos.

Em \textbf{segundo lugar}, a análise da regulação da IA concentrou-se em instrumentos normativos já existentes ou em tramitação (LGPD, MCI, PL 2.338/2023), sem aprofundar o exame de normas técnicas, autorregulatórias ou de soft law que também compõem o ecossistema regulatório.

Em \textbf{terceiro lugar}, a pesquisa adotou uma abordagem predominantemente teórica e analítica, sem realizar estudos empíricos sobre a aplicação concreta dos instrumentos normativos examinados. Estudos de caso sobre decisões judiciais ou administrativas envolvendo IA e direitos da personalidade poderiam oferecer uma perspectiva complementar relevante.

Em \textbf{quarto lugar}, a \MH{} é um documento recente (publicado em 25 de maio de 2026), e a literatura secundária sobre ela ainda é incipiente. À medida que a recepção acadêmica e eclesial da Encíclica se desenvolva, novas interpretações e aplicações poderão enriquecer o debate aqui iniciado.

\section{Agendas de Pesquisa Futura}
\label{sec:agendas-futuras}

Os resultados deste trabalho abrem diversas possibilidades de investigação futura, que podem ser organizadas em quatro eixos.

O \textbf{primeiro eixo} diz respeito ao aprofundamento teórico dos conceitos aqui explorados. A irredutibilidade da pessoa aos dados, o paradigma tecnocrático e os princípios éticos da \MH{} merecem desenvolvimento sistemático em diálogo com a filosofia do direito, a teoria dos direitos fundamentais e a ética aplicada.

O \textbf{segundo eixo} concerne à análise empírica da aplicação dos instrumentos jurídicos examinados. Estudos de caso sobre decisões judiciais envolvendo discriminação algorítmica, direito à explicação e responsabilidade civil por danos causados por sistemas de IA poderiam oferecer uma perspectiva complementar à análise teórica aqui empreendida.

O \textbf{terceiro eixo} refere-se ao direito comparado. A investigação sobre como outros ordenamentos jurídicos (especialmente o europeu (AI Act), o norte-americano e o chinês) têm enfrentado os desafios da IA à personalidade humana poderia oferecer subsídios valiosos para o aperfeiçoamento da regulação brasileira.

O \textbf{quarto eixo} diz respeito ao desenvolvimento tecnológico. À medida que novas aplicações de IA surjam e que os sistemas existentes se tornem mais sofisticados, novas questões de proteção da personalidade certamente emergirão, exigindo permanente atualização do marco regulatório e da reflexão ética que o fundamenta.

\section{Considerações Finais}
\label{sec:consideracoes-finais}

A inteligência artificial não é, em si mesma, uma ameaça à personalidade humana. Como recorda a \MH{} (§89), a tecnologia é uma expressão da criatividade humana e pode servir ao desenvolvimento integral da pessoa quando orientada por valores éticos adequados. O problema não está na tecnologia em si, mas na visão de mundo que a informa e nos usos que dela se fazem.

A \MH{} oferece uma contribuição singular para este debate não porque propõe soluções técnicas ou regras específicas, mas porque recorda o fundamento último de toda proteção jurídica da pessoa: a dignidade que lhe é inerente, anterior a qualquer reconhecimento estatal, independente de qualquer classificação algorítmica.

O desafio que se coloca aos juristas, legisladores e operadores do direito é o de traduzir essa fundamentação antropológica em instrumentos normativos eficazes, capazes de proteger a pessoa humana na era da inteligência artificial. Este trabalho espera ter oferecido uma contribuição para esse esforço coletivo, cuja urgência e importância só tendem a crescer à medida que a IA se integra cada vez mais profundamente em todos os âmbitos da vida social.


% ============================================================
% REFERÊNCIAS
% ============================================================
\bibliography{refs}

% ============================================================
% ANEXOS
% ============================================================
\begin{anexosenv}
    \partanexos
    % ============================================================
% ANEXO — FICHAMENTOS SELECIONADOS
% ============================================================

\chapter{Fichamentos Selecionados}
\label{anexo:fichamentos}

Este anexo reúne os fichamentos dos principais textos que fundamentam a presente pesquisa. Cada fichamento segue o protocolo estabelecido em \texttt{fichamentos/README.md}, contendo a referência completa, os conceitos-chave extraídos, as passagens mais relevantes e a articulação com o problema de pesquisa.

\section{Magistério Pontifício}

\begin{fichamento}
    \textbf{Ref.:} \cite{LeaoXIV2026MH}
    
    \textbf{Conceitos-chave:} dignidade humana, paradigma tecnocrático, irredutibilidade da pessoa, desenvolvimento integral, ecologia integral, destino universal dos bens, subsidiariedade, solidariedade.
    
    \textbf{Passagens relevantes:}
    \begin{itemize}
        \item §3: ``A pessoa humana [...] não pode ser reduzida à soma de seus dados, a um perfil calculável ou a um conjunto de preferências observáveis.''
        \item §89: ``A tecnologia, expressão da criatividade humana, não é em si mesma boa nem má; tudo depende da orientação ética que lhe é dada.''
        \item §101: ``A irredutibilidade da pessoa é o primeiro princípio de toda regulação que pretenda proteger a dignidade humana na era digital.''
        \item §186: ``A pessoa não pode ser transformada em objeto de classificação algorítmica.''
        \item §232: ``Toda a técnica deve estar a serviço do desenvolvimento integral da pessoa humana.''
    \end{itemize}
    
    \textbf{Articulação:} Documento central da pesquisa. Oferece o referencial teórico e principiológico para a análise da proteção da personalidade na era da IA.
\end{fichamento}

\begin{fichamento}
    \textbf{Ref.:} \cite{Francisco2015LS}
    
    \textbf{Conceitos-chave:} ecologia integral, paradigma tecnocrático, cuidado da casa comum, destino universal dos bens, diálogo intergeracional.
    
    \textbf{Passagens relevantes:}
    \begin{itemize}
        \item §106–114: Crítica ao paradigma tecnocrático como visão de mundo que reduz a realidade a objeto de manipulação.
        \item §115: Identificação da raiz humana da crise ecológica.
        \item §203–204: Exigência de uma ecologia integral que articule a dimensão ambiental, social e humana.
    \end{itemize}
    
    \textbf{Articulação:} A crítica ao paradigma tecnocrático presente na \textit{Laudato Si'} é retomada e atualizada pela \MH{} no contexto específico da IA. Constitui a base teórica para a análise crítica do capítulo quinto.
\end{fichamento}

\begin{fichamento}
    \textbf{Ref.:} \cite{JoaoXXIII1963PT}
    
    \textbf{Conceitos-chave:} bem comum, dignidade humana, autoridade política, participação, subsidiariedade.
    
    \textbf{Passagens relevantes:}
    \begin{itemize}
        \item §8–9: Definição do bem comum como conjunto de condições sociais que permitem o desenvolvimento integral da pessoa.
        \item §51–52: A autoridade política como serviço ao bem comum.
        \item §64–65: O princípio da subsidiariedade como limite à atuação do Estado.
    \end{itemize}
    
    \textbf{Articulação:} A doutrina social da Igreja sobre o bem comum e a subsidiariedade, desenvolvida na \textit{Pacem in Terris}, fundamenta os princípios éticos que a \MH{} aplica à regulação da IA.
\end{fichamento}

\begin{fichamento}
    \textbf{Ref.:} \cite{JoaoPauloII1991CA}
    
    \textbf{Conceitos-chave:} livre iniciativa, destino universal dos bens, subsidiariedade, solidariedade, Estado regulador.
    
    \textbf{Passagens relevantes:}
    \begin{itemize}
        \item §15: Crítica ao capitalismo desregulado.
        \item §48: O Estado como promotor do desenvolvimento, não mero espectador.
    \end{itemize}
    
    \textbf{Articulação:} Os princípios da subsidiariedade e da solidariedade, articulados na \textit{Centesimus Annus}, são retomados pela \MH{} como critérios para a governança algorítmica.
\end{fichamento}

\begin{fichamento}
    \textbf{Ref.:} \cite{Francisco2020FT}
    
    \textbf{Conceitos-chave:} fraternidade universal, amizade social, diálogo intergeracional, bem comum, cultura do encontro, solidariedade.
    
    \textbf{Passagens relevantes:}
    \begin{itemize}
        \item §20–21: ``O sonho de fraternidade'' que fundamenta a convivência humana para além de qualquer fronteira.
        \item §34–36: O amor universal como princípio ético que reconhece cada pessoa como irmã, independentemente de sua origem ou condição.
        \item §55: A regra de ouro — não fazer ao outro o que não queres que te façam — como princípio fundamental do diálogo intercultural.
        \item §101–104: A amizade social como condição para a construção de um mundo mais justo e solidário.
        \item §205: Crítica ao individualismo contemporâneo e apelo à solidariedade como virtude política fundamental.
    \end{itemize}
    
    \textbf{Articulação:} A fraternidade universal proposta pela \textit{Fratelli Tutti} oferece um princípio ético fundamental para orientar a governança da IA, assegurando que o desenvolvimento tecnológico esteja a serviço da solidariedade e do bem comum, conforme retomado pela \MH{} no contexto digital.
\end{fichamento}

\begin{fichamento}
    \textbf{Ref.:} \cite{ConcilioVaticanoII1965GS}
    
    \textbf{Conceitos-chave:} dignidade da pessoa humana, progresso técnico, bem comum, autonomia das realidades terrenas, relação Igreja-mundo, desenvolvimento integral.
    
    \textbf{Passagens relevantes:}
    \begin{itemize}
        \item §26: O bem comum como ``conjunto das condições da vida social que permitem aos grupos e a cada um dos seus membros atingir de modo mais completo e facilmente a própria perfeição.''
        \item §34: O progresso técnico como expressão da criatividade humana, que deve estar ordenado ao serviço da pessoa.
        \item §36: A autonomia legítima das realidades terrenas — a ciência e a técnica têm suas leis próprias, mas devem submeter-se aos valores morais fundamentais.
        \item §63: As transformações sociais decorrentes do progresso técnico e seus desafios éticos.
        \item §76: A autonomia da ordem temporal e sua relação com a missão da Igreja na promoção da dignidade humana.
    \end{itemize}
    
    \textbf{Articulação:} A Constituição Pastoral \textit{Gaudium et Spes} estabelece o princípio da autonomia legítima das realidades terrenas (§36) que a \MH{} retoma ao reconhecer o valor intrínseco da técnica, mas também a necessidade de subordiná-la aos imperativos éticos fundamentais, especialmente a dignidade da pessoa humana diante do reducionismo algorítmico.
\end{fichamento}

\section{Livros e Capítulos de Livros}

\begin{fichamento}
    \textbf{Ref.:} \cite{Bioni2021PPD}
    
    \textbf{Conceitos-chave:} proteção de dados pessoais, autodeterminação informativa, privacidade, LGPD, direito à explicação.
    
    \textbf{Conceitos principais:}
    \begin{itemize}
        \item A autodeterminação informativa como fundamento da proteção de dados.
        \item O direito à explicação como garantia do titular diante de decisões automatizadas.
        \item A LGPD como marco regulatório brasileiro da proteção de dados.
    \end{itemize}
    
    \textbf{Articulação:} Fundamenta a análise crítica da LGPD no capítulo quarto e a articulação com os princípios da \MH{} no capítulo quinto.
\end{fichamento}

\begin{fichamento}
    \textbf{Ref.:} \cite{Doneda2006PDD}
    
    \textbf{Conceitos-chave:} proteção de dados, privacidade, direito comparado, evolução legislativa.
    
    \textbf{Conceitos principais:}
    \begin{itemize}
        \item A proteção de dados como direito autônomo.
        \item A evolução legislativa europeia e brasileira.
        \item A necessidade de harmonização entre proteção de dados e direitos da personalidade.
    \end{itemize}
    
    \textbf{Articulação:} O marco teórico da proteção de dados desenvolvido por Doneda fornece as bases conceituais para a análise do regime jurídico brasileiro no capítulo quarto.
\end{fichamento}

\begin{fichamento}
    \textbf{Ref.:} \cite{Rodota2008PDA}
    
    \textbf{Conceitos-chave:} privacidade, proteção de dados, dignidade, identidade pessoal, sociedade informacional.
    
    \textbf{Conceitos principais:}
    \begin{itemize}
        \item A privacidade como condição para o exercício da liberdade.
        \item A proteção de dados como instrumento de tutela da identidade pessoal.
        \item O direito de controlar a própria informação como dimensão da dignidade.
    \end{itemize}
    
    \textbf{Articulação:} A conexão entre proteção de dados e dignidade humana, sustentada por Rodotà, dialoga diretamente com a fundamentação antropológica da \MH{}, especialmente a tese da irredutibilidade da pessoa aos seus dados.
\end{fichamento}

\begin{fichamento}
    \textbf{Ref.:} \cite{Castells2010EI}
    
    \textbf{Conceitos-chave:} sociedade em rede, era da informação, poder, identidade, globalização.
    
    \textbf{Conceitos principais:}
    \begin{itemize}
        \item A sociedade em rede como nova estrutura social da era da informação.
        \item O poder dos fluxos de informação como nova forma de dominação.
        \item A crise da identidade pessoal na sociedade em rede.
    \end{itemize}
    
    \textbf{Articulação:} O quadro sociológico de Castells fornece o contexto para compreender os desafios à personalidade humana na era digital, examinados no capítulo terceiro.
\end{fichamento}

\begin{fichamento}
    \textbf{Ref.:} \cite{Zuboff2019SC}
    
    \textbf{Conceitos-chave:} capitalismo de vigilância, excedente comportamental, sociedade instrumental, direitos humanos.
    
    \textbf{Conceitos principais:}
    \begin{itemize}
        \item O capitalismo de vigilância como nova ordem econômica que trata a experiência humana como matéria-prima gratuita.
        \item A instrumentalização da pessoa como ameaça aos direitos humanos.
        \item A necessidade de um novo contrato social para a era digital.
    \end{itemize}
    
    \textbf{Articulação:} A teoria do capitalismo de vigilância fundamenta a análise dos riscos da IA à personalidade no capítulo terceiro e dialoga com a crítica da \MH{} ao paradigma tecnocrático no capítulo quinto.
\end{fichamento}

\begin{fichamento}
    \textbf{Ref.:} \cite{AquinoST}
    
    \textbf{Conceitos-chave:} pessoa, dignidade, lei natural, bem comum, substância individual de natureza racional, sindérese, virtude.
    
    \textbf{Conceitos principais:}
    \begin{itemize}
        \item A definição de pessoa como ``substância individual de natureza racional'' (\textit{individua substantia rationalis naturae}), que estabelece a dignidade ontológica do ser humano (I, q. 29, a. 1).
        \item A lei natural como participação da lei eterna na criatura racional, fundamento objetivo da moralidade (I-II, q. 91, a. 2).
        \item O bem comum como fim da sociedade política e critério de justiça das leis positivas (I-II, q. 90, a. 2).
        \item A sindérese como hábito natural dos primeiros princípios práticos, base da consciência moral (I, q. 79, a. 12).
        \item A virtude como aperfeiçoamento das faculdades humanas para o agir segundo a razão reta (I-II, q. 55, a. 1).
    \end{itemize}
    
    \textbf{Articulação:} A antropologia filosófica de Tomás de Aquino, em particular a definição de pessoa como substância individual de natureza racional, fundamenta a tese central da \MH{} sobre a irredutibilidade da pessoa face aos sistemas de IA. O conceito de lei natural informa a crítica ao paradigma tecnocrático, que subordina o bem comum à lógica algorítmica, enquanto a doutrina do bem comum fornece o critério teleológico para a regulação ética da inteligência artificial.
\end{fichamento}

\begin{fichamento}
    \textbf{Ref.:} \cite{Mounier1949P}
    
    \textbf{Conceitos-chave:} pessoa, engajamento, comunidade, liberdade, personalismo comunitário, transcendência, vocação.
    
    \textbf{Conceitos principais:}
    \begin{itemize}
        \item A pessoa como ``ser espiritual, naturalmente independente, transcendente às realidades materiais e sociais, mas que se realiza na comunhão com as outras pessoas''.
        \item A distinção entre individualidade (fechamento egoísta) e personalidade (abertura ao outro e à comunidade).
        \item O engajamento pessoal como condição da liberdade autêntica — a pessoa se realiza na ação responsável.
        \item A comunidade como meio e fim do desenvolvimento pessoal, contraposta ao individualismo liberal e ao coletivismo totalitário.
        \item A transcendência como dimensão constitutiva da pessoa que a abre ao absoluto e ao infinito.
    \end{itemize}
    
    \textbf{Articulação:} O personalismo comunitário de Mounier fornece o quadro filosófico que informa a concepção de pessoa na \MH{}. A crítica mounieriana à redução da pessoa a funções sociais ou econômicas antecipa a crítica da encíclica à redução algorítmica da pessoa a perfis de dados. A comunidade como realização da pessoa orienta a visão de uma governança algorítmica participativa e solidária.
\end{fichamento}

\begin{fichamento}
    \textbf{Ref.:} \cite{Ricoeur1990SM}
    
    \textbf{Conceitos-chave:} identidade narrativa, ipseidade, alteridade, estima de si, solicitude, sabedoria prática.
    
    \textbf{Conceitos principais:}
    \begin{itemize}
        \item A distinção entre identidade-idem (\textit{idem} — mesmidade, caráter fixo) e identidade-ipse (\textit{ipse} — ipseidade, si-mesmo constituído na promessa e na narrativa).
        \item A ipseidade como identidade dinâmica que se constitui na relação com o outro e na fidelidade à palavra dada.
        \item A alteridade constitutiva do si-mesmo — ``si-mesmo como um outro'' —, que insere a dimensão ética no núcleo da identidade pessoal.
        \item A estima de si e a solicitude como mediações éticas da identidade narrativa.
        \item A sabedoria prática como julgamento em situação concreta, para além da aplicação mecânica de regras universais.
    \end{itemize}
    
    \textbf{Articulação:} A noção de identidade narrativa de Ricoeur fundamenta a crítica à redução da pessoa a perfis algorítmicos na era da IA. Se a identidade pessoal se constitui narrativamente, através de um processo contínuo de interpretação e reinterpretação, então os sistemas de IA que classificam, predizem e definem a pessoa com base em dados passados cometem um erro antropológico fundamental: confundem a mesmidade (identidade-idem) com a totalidade da pessoa, suprimindo a dimensão de abertura e novidade própria da ipseidade.
\end{fichamento}

\begin{fichamento}
    \textbf{Ref.:} \cite{Jonas1979PR}
    
    \textbf{Conceitos-chave:} responsabilidade, ética do futuro, imperativo da responsabilidade, heurística do medo, vulnerabilidade, precaução.
    
    \textbf{Conceitos principais:}
    \begin{itemize}
        \item O imperativo da responsabilidade: ``Age de tal modo que os efeitos da tua ação sejam compatíveis com a permanência de uma vida humana autêntica sobre a Terra.''
        \item A heurística do medo (\textit{Heuristik der Furcht}) como método para antecipar cenários catastróficos antes que se concretizem.
        \item A vulnerabilidade da natureza e das gerações futuras como fundamento de uma nova ética da responsabilidade.
        \item A crítica ao utopismo tecnológico que ignora os riscos de longo prazo do desenvolvimento técnico.
        \item A assimetria entre o poder da técnica e a capacidade humana de prever suas consequências.
    \end{itemize}
    
    \textbf{Articulação:} A ética da responsabilidade de Jonas fundamenta o princípio de precaução aplicado à IA, retomado pela \MH{} na crítica ao paradigma tecnocrático. A heurística do medo oferece um método para identificar riscos existenciais da IA antes que se materializem, enquanto o imperativo da responsabilidade estabelece o dever ético de assegurar que o desenvolvimento algorítmico seja compatível com a dignidade e a autonomia da pessoa humana, orientando a regulação prudencial proposta no capítulo quinto.
\end{fichamento}

\begin{fichamento}
    \textbf{Ref.:} \cite{Floridi2023EI}
    
    \textbf{Conceitos-chave:} ética da IA, princípios éticos, governança algorítmica, dignidade informacional, sociedade da informação, infraética.
    
    \textbf{Conceitos principais:}
    \begin{itemize}
        \item A ética da IA como campo disciplinar autônomo que articula princípios normativos para o desenvolvimento e uso de sistemas inteligentes.
        \item A distinção entre ética aplicada (\textit{applied ethics}) e infraética (\textit{infraethics}), que fornece o enquadramento normativo de fundo para a governança algorítmica.
        \item A dignidade informacional como princípio ético fundamental que reconhece o valor intrínseco de cada pessoa na sociedade da informação.
        \item Os quatro princípios éticos para a IA: beneficência, não maleficência, autonomia e justiça, acrescidos da explicabilidade como princípio transversal.
        \item A necessidade de traduzir princípios éticos em práticas regulatórias efetivas, superando a mera declaração de intenções.
    \end{itemize}
    
    \textbf{Articulação:} Oferece o marco ético contemporâneo da IA que dialoga com os princípios propostos pela \MH{}. A dignidade informacional de Floridi converge com a irredutibilidade da pessoa sustentada pela encíclica, enquanto os princípios éticos propostos pelo autor fornecem um quadro comparativo para os princípios da \MH{} examinados no capítulo segundo, permitindo uma análise crítica dos limites e potencialidades de cada abordagem.
\end{fichamento}

\begin{fichamento}
    \textbf{Ref.:} \cite{MendesDoneda2021PD}
    
    \textbf{Conceitos-chave:} proteção de dados, LGPD, direito comparado, dignidade, autodeterminação informativa, privacidade, ANPD.
    
    \textbf{Conceitos principais:}
    \begin{itemize}
        \item A proteção de dados pessoais como direito fundamental autônomo, para além da mera tutela da privacidade.
        \item A análise crítica da LGPD: avanços e insuficiências do marco regulatório brasileiro diante dos desafios tecnológicos.
        \item O direito à explicação como garantia específica diante de decisões automatizadas e seus limites na legislação brasileira.
        \item A Autoridade Nacional de Proteção de Dados (ANPD) como instituição central para a efetividade do sistema de proteção.
        \item A perspectiva comparada com o direito europeu e norte-americano como método para identificar lacunas e boas práticas.
    \end{itemize}
    
    \textbf{Articulação:} A perspectiva crítica sobre a proteção de dados informa a análise das insuficiências do regime brasileiro diante dos desafios da IA, desenvolvida no capítulo quarto. A ênfase na proteção de dados como direito fundamental autônomo dialoga com a fundamentação antropológica da \MH{}, que reconhece na dignidade da pessoa o fundamento último de toda proteção jurídica na era digital.
\end{fichamento}

\begin{fichamento}
    \textbf{Ref.:} \cite{Comparato2016DH}
    
    \textbf{Conceitos-chave:} direitos humanos, dignidade, evolução histórica, universalidade, indivisibilidade, fundamentação filosófica.
    
    \textbf{Conceitos principais:}
    \begin{itemize}
        \item A afirmação histórica dos direitos humanos como processo de positivação da dignidade da pessoa humana ao longo dos séculos.
        \item A universalidade dos direitos humanos como expressão da igual dignidade de todos os seres humanos, independentemente de sua condição.
        \item A indivisibilidade dos direitos humanos — direitos civis, políticos, econômicos, sociais e culturais formam um conjunto unitário.
        \item A fundamentação filosófica dos direitos humanos na tradição do direito natural e no personalismo ético.
        \item O reconhecimento dos novos direitos diante das transformações técnicas e sociais como característica dinâmica do sistema de direitos humanos.
    \end{itemize}
    
    \textbf{Articulação:} A fundamentação histórica dos direitos humanos desenvolvida por Comparato fornece o pano de fundo para compreender a proteção da personalidade como direito fundamental em evolução. A tese do autor sobre a abertura do sistema de direitos humanos a novas exigências éticas fundamenta a proposta de extensão da tutela da personalidade diante dos desafios específicos da IA, em consonância com a visão da \MH{} sobre o desenvolvimento integral.
\end{fichamento}

\begin{fichamento}
    \textbf{Ref.:} \cite{Sarlet2021DF}
    
    \textbf{Conceitos-chave:} direitos fundamentais, dimensão objetiva, eficácia horizontal, dignidade, proporcionalidade, restrição a direitos.
    
    \textbf{Conceitos principais:}
    \begin{itemize}
        \item A dupla dimensão dos direitos fundamentais: subjetiva (direitos subjetivos) e objetiva (valores objetivos que informam todo o ordenamento).
        \item A eficácia horizontal dos direitos fundamentais como vinculação dos particulares, especialmente relevante para as relações entre plataformas digitais e usuários.
        \item A dignidade da pessoa humana como fundamento e limite de todos os direitos fundamentais, com eficácia irradiante sobre todo o ordenamento jurídico.
        \item O princípio da proporcionalidade como método de solução de conflitos entre direitos fundamentais na era digital.
        \item A proteção contra automações — o direito a não ser submetido a decisões exclusivamente automatizadas — como nova dimensão dos direitos fundamentais.
    \end{itemize}
    
    \textbf{Articulação:} A dogmática dos direitos fundamentais informa a análise da eficácia dos direitos da personalidade diante dos desafios da IA, especialmente no que tange à eficácia horizontal (vinculação de plataformas privadas) e à dupla dimensão objetiva e subjetiva. Esta análise fundamenta a tese, desenvolvida nos capítulos quarto e quinto, de que a proteção da personalidade na era da IA exige tanto mecanismos subjetivos de tutela quanto a adoção de políticas públicas orientadas pelos valores objetivos da ordem constitucional.
\end{fichamento}

\begin{fichamento}
    \textbf{Ref.:} \cite{Frazao2020RD}
    
    \textbf{Conceitos-chave:} digitalização, riscos, regulação, responsabilidade civil, plataformas, assimetria de poder, proteção do consumidor.
    
    \textbf{Conceitos principais:}
    \begin{itemize}
        \item A digitalização como fenômeno jurídico multifacetado que impõe novos riscos ao direito privado e à proteção da pessoa.
        \item A assimetria de poder entre plataformas digitais e usuários como desafio central para a regulação do ambiente digital.
        \item A responsabilidade civil dos provedores de serviços algorítmicos diante de danos à personalidade.
        \item A insuficiência dos modelos tradicionais de responsabilidade baseados na culpa diante dos riscos sistêmicos da IA.
        \item A necessidade de uma abordagem regulatória baseada em risco, que considere o potencial danoso dos sistemas algorítmicos.
    \end{itemize}
    
    \textbf{Articulação:} A análise dos riscos da digitalização desenvolvida por Frazão fundamenta a identificação dos desafios à proteção da personalidade na era da IA, examinados no capítulo terceiro. Sua proposta de regulação baseada em risco dialoga diretamente com a abordagem da \MH{}, que também propõe uma avaliação ética prévia dos sistemas tecnológicos como condição para sua adoção legítima.
\end{fichamento}

\begin{fichamento}
    \textbf{Ref.:} \cite{Pasquale2015BM}
    
    \textbf{Conceitos-chave:} opacidade algorítmica, caixa-preta, assimetria de poder, regulação, accountability, sociedade da informação.
    
    \textbf{Conceitos principais:}
    \begin{itemize}
        \item A opacidade dos algoritmos como problema de governança democrática — decisões que afetam a vida das pessoas são tomadas por sistemas incompreensíveis.
        \item A assimetria de poder entre as corporações tecnológicas e os cidadãos na era dos dados massivos.
        \item A necessidade de transparência algorítmica como condição para a accountability e o controle social dos sistemas automatizados.
        \item O ciclo vicioso entre coleta de dados, profiling e discriminação como ameaça sistêmica à autonomia pessoal.
        \item A proposta de uma ``sociedade da caixa-preta aberta'' (\textit{open black box society}) como ideal regulatório.
    \end{itemize}
    
    \textbf{Articulação:} A crítica à opacidade dos sistemas algorítmicos desenvolvida por Pasquale dialoga com a exigência de transparência da \MH{} para a governança da IA. O conceito de ``sociedade da caixa-preta'' fundamenta a análise crítica do capítulo terceiro sobre os riscos à personalidade, enquanto a proposta de abertura dos sistemas algorítmicos informa as recomendações regulatórias do capítulo quinto, em consonância com o princípio de subsidiariedade defendido pela encíclica.
\end{fichamento}

\begin{fichamento}
    \textbf{Ref.:} \cite{Amaral2003DBC}
    
    \textbf{Conceitos-chave:} pessoa, capacidade, personalidade, direitos da personalidade, bem comum, sujeito de direito.
    
    \textbf{Conceitos principais:}
    \begin{itemize}
        \item A pessoa natural como sujeito de direito — capacidade jurídica como atributo de todo ser humano.
        \item Os direitos da personalidade como direitos inatos, absolutos, intransmissíveis, irrenunciáveis e imprescritíveis.
        \item A distinção entre capacidade de direito (aptidão para ser titular de direitos) e capacidade de fato (aptidão para exercê-los pessoalmente).
        \item O nome, a imagem, a honra, a privacidade e a integridade moral como projeções da personalidade tuteladas pelo direito civil.
        \item A teoria geral dos direitos da personalidade como construção dogmática aberta a novas manifestações da personalidade.
    \end{itemize}
    
    \textbf{Articulação:} A dogmática civilística dos direitos da personalidade, sistematizada por Amaral, fornece o quadro jurídico fundamental para a pesquisa. A abertura do sistema de direitos da personalidade a novas manifestações — como o direito à autodeterminação informativa e à explicação algorítmica — fundamenta a tese, defendida no capítulo quarto, de que a proteção contra os riscos da IA constitui uma nova dimensão dos direitos da personalidade, compatível com a visão integral da pessoa humana proposta pela \MH{}.
\end{fichamento}

\begin{fichamento}
    \textbf{Ref.:} \cite{Noble2018AFS}
    
    \textbf{Conceitos-chave:} discriminação algorítmica, racismo algorítmico, vieses, busca online, desigualdade, justiça social.
    
    \textbf{Conceitos principais:}
    \begin{itemize}
        \item A discriminação algorítmica como reprodução e amplificação das desigualdades sociais existentes através de sistemas automatizados.
        \item O racismo algorítmico — como os sistemas de busca e recomendação perpetuam estereótipos raciais e de gênero.
        \item A opacidade dos critérios de ranqueamento e recomendação como obstáculo à identificação e correção de vieses discriminatórios.
        \item A responsabilidade social das empresas de tecnologia na construção de sistemas equitativos.
        \item A necessidade de intervenção regulatória para assegurar que os sistemas algorítmicos respeitem princípios de justiça e não discriminação.
    \end{itemize}
    
    \textbf{Articulação:} A análise das discriminações algorítmicas desenvolvida por Noble fundamenta a necessidade de regulação da IA que a \MH{} defende. O estudo dos vieses algorítmicos na busca online ilustra uma das ameaças concretas à personalidade examinadas no capítulo terceiro — a objetificação algorítmica que reduz a pessoa a estereótipos — e reforça o argumento da encíclica sobre a necessidade de submeter a técnica a critérios éticos que respeitem a igual dignidade de todos os seres humanos.
\end{fichamento}

\begin{fichamento}
    \textbf{Ref.:} \cite{Pariser2022THP}
    
    \textbf{Conceitos-chave:} bolha de filtros, personalização, manipulação, autonomia, economia da atenção, fragmentação social.
    
    \textbf{Conceitos principais:}
    \begin{itemize}
        \item A bolha de filtros (\textit{filter bubble}) como isolamento informacional causado pela personalização algorítmica do conteúdo online.
        \item A personalização algorítmica como ameaça à autonomia pessoal — o usuário não escolhe o que vê, mas é conduzido por critérios opacos.
        \item A fragmentação do espaço público em bolhas incomunicantes como risco à democracia e ao diálogo social.
        \item A economia da atenção como motor dos sistemas de personalização — o conteúdo é selecionado para maximizar o engajamento, não para informar ou formar.
        \item A perda de controle sobre o próprio ambiente informacional como dimensão da crise da autonomia na era digital.
    \end{itemize}
    
    \textbf{Articulação:} O fenômeno das bolhas informacionais identificado por Pariser ilustra a ameaça à autonomia pessoal que a \MH{} denuncia no paradigma digital. A fragmentação do espaço público pela personalização algorítmica contradiz a vocação comunitária e dialógica da pessoa humana afirmada pela encíclica, enquanto a crítica à economia da atenção fundamenta a necessidade de uma regulação ética dos sistemas algorítmicos de recomendação.
\end{fichamento}

\begin{fichamento}
    \textbf{Ref.:} \cite{Kurzweil2005THS}
    
    \textbf{Conceitos-chave:} singularidade tecnológica, transumanismo, superinteligência, imortalidade, convergência NBIC, pós-humanismo.
    
    \textbf{Conceitos principais:}
    \begin{itemize}
        \item A lei dos retornos acelerados — o progresso tecnológico segue uma trajetória exponencial que levará à singularidade por volta de 2045.
        \item A singularidade tecnológica como momento em que a inteligência artificial superará a inteligência humana em todas as dimensões.
        \item O transumanismo como projeto de superação dos limites biológicos humanos através da fusão com máquinas inteligentes.
        \item A imortalidade técnica como objetivo último do projeto transumanista — o upload da mente humana para suportes digitais.
        \item A convergência NBIC (nanotecnologia, biotecnologia, tecnologia da informação e ciência cognitiva) como base material da transformação pós-humana.
    \end{itemize}
    
    \textbf{Articulação:} O transumanismo de Kurzweil representa a visão que a \MH{} critica como reducionismo tecnológico da pessoa humana. A tese da singularidade tecnológica e da superação dos limites biológicos contradiz a antropologia integral da encíclica, que afirma a irredutibilidade da pessoa à sua dimensão meramente biológica ou informacional. A análise do transumanismo no capítulo terceiro permite demonstrar que o projeto de fusão homem-máquina representa a forma mais radical de objetificação algorítmica, contra a qual a \MH{} oferece uma alternativa fundada na dignidade transcendente da pessoa.
\end{fichamento}

\begin{fichamento}
    \textbf{Ref.:} \cite{Sarmento2016DPF}
    
    \textbf{Conceitos-chave:} dignidade, mínimo existencial, autonomia, reconhecimento, vulnerabilidade, igualdade.
    
    \textbf{Conceitos principais:}
    \begin{itemize}
        \item A dignidade da pessoa humana como princípio constitucional fundamental que irradia efeitos sobre todo o ordenamento jurídico.
        \item A dupla dimensão da dignidade: negativa (não instrumentalização) e positiva (promoção das condições de vida digna).
        \item O mínimo existencial como núcleo intangível da dignidade que inclui tanto condições materiais quanto elementos imateriais (reconhecimento, autonomia, integridade moral).
        \item A autonomia pessoal como expressão da dignidade — capacidade de autodeterminação que deve ser protegida inclusive contra ingerências tecnológicas.
        \item O reconhecimento intersubjetivo como dimensão essencial da dignidade — a pessoa se realiza na relação com o outro e precisa ser reconhecida em seu valor intrínseco.
    \end{itemize}
    
    \textbf{Articulação:} A fundamentação da dignidade como direito fundamental desenvolvida por Sarmento informa a crítica da \MH{} à instrumentalização algorítmica da pessoa. A dupla dimensão da dignidade (negativa e positiva) fornece o quadro para compreender tanto as violações diretas da personalidade pela IA (vigilância, discriminação) quanto as omissões estatais na promoção de condições de autodeterminação digital. O reconhecimento como elemento da dignidade fundamenta a exigência de que os sistemas algorítmicos tratem cada pessoa como fim, jamais como meio.
\end{fichamento}

\begin{fichamento}
    \textbf{Ref.:} \cite{Doneda2022DPD}
    
    \textbf{Conceitos-chave:} direitos da personalidade, sociedade da informação, proteção de dados, autonomia privada, identidade pessoal, consentimento.
    
    \textbf{Conceitos principais:}
    \begin{itemize}
        \item A evolução dos direitos da personalidade na sociedade informacional — do direito à privacidade ao direito à autodeterminação informativa.
        \item O consentimento como instrumento de tutela insuficiente diante das assimetrias de poder e informação no ambiente digital.
        \item A identidade pessoal como construção social ameaçada pela redução a perfis algorítmicos.
        \item A necessidade de repensar os direitos da personalidade à luz das novas tecnologias, sem perder sua fundamentação antropológica.
        \item O direito ao esquecimento e o direito à explicação como novas manifestações dos direitos da personalidade na era digital.
    \end{itemize}
    
    \textbf{Articulação:} A análise dos direitos da personalidade na sociedade informacional desenvolvida por Doneda fornece o quadro teórico para compreender os desafios regulatórios da IA examinados no capítulo quarto. Sua crítica ao consentimento como instrumento insuficiente de tutela dialoga com a necessidade, apontada pela \MH{}, de uma regulação ética que vá além da mera autonomia individual e considere o bem comum como critério orientador do desenvolvimento tecnológico.
\end{fichamento}

\begin{fichamento}
    \textbf{Ref.:} \cite{Maritain1947PCG}
    
    \textbf{Conceitos-chave:} pessoa, bem comum, dignidade, sociabilidade, individualidade, personalidade, transcendência.
    
    \textbf{Conceitos principais:}
    \begin{itemize}
        \item A distinção entre individualidade (aquilo que no homem é matéria, princípio de individuação e egoísmo) e personalidade (aquilo que no homem é espírito, abertura ao outro e ao absoluto).
        \item A pessoa como ``intelecto e vontade'', ser espiritual que transcende a ordem material e social.
        \item O bem comum como fim da sociedade que não se reduz à soma dos bens individuais, mas inclui a realização integral de cada pessoa.
        \item A sociabilidade constitutiva da pessoa — o homem é naturalmente social e se realiza na comunidade política.
        \item A subordinação do bem particular ao bem comum e a simultânea transcendência da pessoa em relação à sociedade política.
    \end{itemize}
    
    \textbf{Articulação:} A distinção entre individualidade e personalidade em Maritain fundamenta a tese da irredutibilidade da pessoa que a \MH{} sustenta diante do reducionismo algorítmico. Se a personalidade é abertura espiritual e transcendência, então qualquer sistema de IA que pretenda capturar a totalidade da pessoa através de dados e algoritmos comete um erro ontológico fundamental. O bem comum como fim da sociedade orienta a concepção de governança algorítmica proposta no capítulo quinto.
\end{fichamento}

\begin{fichamento}
    \textbf{Ref.:} \cite{Wojtyla1979PA}
    
    \textbf{Conceitos-chave:} pessoa, ação, consciência, transcendência, dignidade, autodeterminação, experiência vivida.
    
    \textbf{Conceitos principais:}
    \begin{itemize}
        \item A pessoa se manifesta através da ação — a ação é o meio privilegiado para conhecer a pessoa e sua estrutura ontológica.
        \item A transcendência da pessoa na ação — o ser humano não é apenas determinado por causas externas, mas se autodetermina através de escolhas livres e responsáveis.
        \item A consciência como elemento constitutivo da ação pessoal — a pessoa age com conhecimento e vontade, diferentemente de um mero processador de inputs.
        \item A autodeterminação como capacidade de agir sobre si mesmo, de se autogovernar, de ser causa de suas próprias ações.
        \item A dignidade da pessoa como sujeito moral — a pessoa é fins em si mesma, jamais meio para fins alheios.
    \end{itemize}
    
    \textbf{Articulação:} A fenomenologia da ação pessoal de Karol Wojtyła informa a compreensão da pessoa como ser livre e autorresponsável, que a \MH{} contrapõe à visão determinista dos algoritmos. A tese da transcendência na ação — o ser humano não se reduz a inputs e outputs — fundamenta a crítica radical da encíclica à objetificação algorítmica, que ignora a capacidade humana de autodeterminação e reduz a pessoa a um conjunto de comportamentos previsíveis e manipuláveis.
\end{fichamento}

\section{Artigos Científicos}

\begin{fichamento}
    \textbf{Ref.:} \cite{Floridi2018AIG}
    
    \textbf{Conceitos-chave:} ética da IA, princípios éticos, beneficência, não maleficência, autonomia, justiça, explicabilidade.
    
    \textbf{Conceitos principais:}
    \begin{itemize}
        \item Quatro princípios éticos para a IA: beneficência, não maleficência, autonomia e justiça.
        \item A explicabilidade como princípio transversal.
        \item A necessidade de traduzir princípios éticos em práticas regulatórias.
    \end{itemize}
    
    \textbf{Articulação:} Oferece um quadro comparativo para os princípios éticos da \MH{}, examinados no capítulo segundo e aplicados no capítulo quinto.
\end{fichamento}

\begin{fichamento}
    \textbf{Ref.:} \cite{Calo2017AIL}
    
    \textbf{Conceitos-chave:} IA, regulação, direito, mercado, segurança, responsabilidade.
    
    \textbf{Conceitos principais:}
    \begin{itemize}
        \item Desafios regulatórios específicos da IA.
        \item A inadequação dos modelos regulatórios tradicionais.
        \item A necessidade de regulação adaptativa.
    \end{itemize}
    
    \textbf{Articulação:} Fundamenta a análise das insuficiências do ordenamento jurídico brasileiro diante dos desafios da IA, no capítulo quarto.
\end{fichamento}

\begin{fichamento}
    \textbf{Ref.:} \cite{Mulholland2021IA}
    
    \textbf{Conceitos-chave:} inteligência artificial, direito, regulação, responsabilidade, direitos fundamentais, dignidade.
    
    \textbf{Conceitos principais:}
    \begin{itemize}
        \item Os desafios específicos da IA para o direito brasileiro: opacidade, imprevisibilidade e autonomia dos sistemas algorítmicos.
        \item A responsabilidade civil por danos causados por sistemas de IA — da responsabilidade subjetiva à responsabilidade objetiva e ao risco criado.
        \item A insuficiência dos institutos tradicionais do direito privado diante dos riscos sistêmicos da inteligência artificial.
        \item A necessidade de uma abordagem interdisciplinar que articule direito, ética e tecnologia para a regulação da IA.
        \item A proteção de dados pessoais como primeiro passo para a regulação da IA no Brasil.
    \end{itemize}
    
    \textbf{Articulação:} O artigo fundamenta a análise dos desafios jurídicos específicos da IA no direito brasileiro, desenvolvida no capítulo quarto. A proposta de uma abordagem interdisciplinar dialoga com a perspectiva integral da \MH{}, que reconhece a complexidade da relação entre técnica, ética e direito na proteção da personalidade humana.
\end{fichamento}

\begin{fichamento}
    \textbf{Ref.:} \cite{Wachter2017TSP}
    
    \textbf{Conceitos-chave:} direito à explicação, GDPR, decisões automatizadas, transparência, opacidade algorítmica, profiling.
    
    \textbf{Conceitos principais:}
    \begin{itemize}
        \item O direito à explicação sob o GDPR: alcance e limites do Art. 22 do Regulamento Europeu.
        \item A distinção entre explicação \textit{ex ante} (antes da decisão) e explicação \textit{ex post} (após a decisão) como graus diversos de proteção do titular.
        \item A insuficiência do direito à explicação nos casos de modelos algorítmicos complexos (\textit{black box}).
        \item A tensão entre o direito à explicação e os segredos comerciais protegidos pelas empresas de tecnologia.
        \item A proposta de um ``direito à informação significativa'' (\textit{right to meaningful information}) sobre o funcionamento dos sistemas algorítmicos.
    \end{itemize}
    
    \textbf{Articulação:} A defesa do direito à explicação como garantia fundamental nas decisões automatizadas fundamenta a análise da LGPD no capítulo quarto e dialoga com o princípio de transparência proposto pela \MH{}. As limitações identificadas pelas autoras informam a necessidade de um aperfeiçoamento regulatório que vá além do mero direito à explicação e assegure uma transparência substancial dos sistemas algorítmicos.
\end{fichamento}

\begin{fichamento}
    \textbf{Ref.:} \cite{Barocas2016BFF}
    
    \textbf{Conceitos-chave:} discriminação, big data, viés algorítmico, equidade, justiça, disparate impact.
    
    \textbf{Conceitos principais:}
    \begin{itemize}
        \item O impacto discriminatório dos sistemas de big data — mesmo sem intenção de discriminar, os algoritmos podem perpetuar e amplificar desigualdades existentes.
        \item A dificuldade de detectar discriminação algorítmica quando os critérios de decisão são opacos e os dados históricos já refletem preconceitos sociais.
        \item A distinção entre tratamento desigual (\textit{disparate treatment}) e impacto desproporcional (\textit{disparate impact}) como categorias analíticas para avaliar a justiça algorítmica.
        \item O ``ciclo de feedback'' entre dados discriminatórios e decisões algorítmicas — o viés se autoalimenta.
        \item A insuficiência das abordagens meramente técnicas para a equidade algorítmica — a justiça algorítmica exige também intervenção normativa.
    \end{itemize}
    
    \textbf{Articulação:} A análise dos vieses discriminatórios em sistemas automatizados desenvolvida por Barocas fundamenta a crítica da \MH{} à objetificação algorítmica da pessoa. O conceito de impacto desproporcional (\textit{disparate impact}) oferece um critério jurídico para avaliar a equidade dos sistemas de IA, complementar ao princípio de justiça defendido pela encíclica, informando as propostas regulatórias do capítulo quinto.
\end{fichamento}

\begin{fichamento}
    \textbf{Ref.:} \cite{Chesney2019DF}
    
    \textbf{Conceitos-chave:} desinformação, deep fakes, regulação, democracia, segurança, manipulação, autenticidade.
    
    \textbf{Conceitos principais:}
    \begin{itemize}
        \item As deep fakes como nova forma de desinformação que utiliza IA para criar conteúdos audiovisuais falsos indistinguíveis da realidade.
        \item Os riscos das deep fakes para a democracia: manipulação eleitoral, destruição de reputações e erosão da confiança pública.
        \item As deep fakes como ameaça à identidade pessoal — a possibilidade de atribuir falsamente atos e palavras a uma pessoa.
        \item Os desafios regulatórios específicos das deep fakes: detecção técnica insuficiente, liberdade de expressão e jurisdição transnacional.
        \item A necessidade de uma resposta combinada de tecnologia, direito e educação midiática para enfrentar o fenômeno.
    \end{itemize}
    
    \textbf{Articulação:} O fenômeno das deep fakes ilustra uma das ameaças concretas à personalidade que a \MH{} aborda indiretamente — a manipulação algorítmica que atinge a identidade, a honra e a imagem da pessoa. A análise dos riscos das deep fakes fundamenta a necessidade, defendida no capítulo quinto, de uma regulação específica que proteja a personalidade contra as formas mais invasivas de manipulação algorítmica.
\end{fichamento}

\begin{fichamento}
    \textbf{Ref.:} \cite{DoshiVelez2017MAR}
    
    \textbf{Conceitos-chave:} accountability, IA, responsabilidade, explicabilidade, prestação de contas, governança.
    
    \textbf{Conceitos principais:}
    \begin{itemize}
        \item A accountability (prestação de contas) como princípio fundamental para a governança dos sistemas de IA.
        \item A tensão entre a complexidade técnica dos modelos de IA e a exigibilidade de accountability pública.
        \item A necessidade de mecanismos de auditoria independente para sistemas algorítmicos que afetam direitos fundamentais.
        \item O direito à explicação como instrumento de accountability — a transparência algorítmica como condição para a responsabilização.
        \item A proposta de standards mínimos de accountability para diferentes níveis de risco dos sistemas de IA.
    \end{itemize}
    
    \textbf{Articulação:} O artigo fundamenta a necessidade de sistemas de IA accountable, compatível com a exigência de responsabilidade (\textit{accountability}) da \MH{}. A proposta de standards mínimos de accountability por nível de risco dialoga diretamente com a classificação de risco proposta pelo PL 2.338/2023 e analisada no capítulo quarto, oferecendo critérios para seu aperfeiçoamento à luz dos princípios éticos da encíclica.
\end{fichamento}

\begin{fichamento}
    \textbf{Ref.:} \cite{DoshiVelez2017TMI}
    
    \textbf{Conceitos-chave:} IA interpretável, explicabilidade, rigor científico, complexidade, transparência, confiança.
    
    \textbf{Conceitos principais:}
    \begin{itemize}
        \item A distinção entre modelos intrinsecamente interpretáveis (como árvores de decisão) e modelos pós-hoc explicáveis (explicações externas a modelos complexos).
        \item A tensão fundamental entre acurácia e interpretabilidade — modelos mais precisos tendem a ser menos interpretáveis.
        \item A necessidade de um programa de pesquisa rigoroso para a IA explicável (\textit{explainable AI — XAI}).
        \item A identificação de diferentes públicos para a explicação (usuários finais, reguladores, desenvolvedores) e suas diferentes necessidades informacionais.
        \item Os limites da explicabilidade — nem toda decisão algorítmica pode ser explicada de forma completa e compreensível.
    \end{itemize}
    
    \textbf{Articulação:} A exigência de IA interpretável fundamenta o direito à explicação analisado na pesquisa. A distinção entre explicabilidade intrínseca e pós-hoc informa a análise crítica dos limites do direito à explicação na LGPD e no PL 2.338/2023, desenvolvida no capítulo quarto. A complexidade identificada pelos autores reforça a necessidade de uma abordagem normativa que a \MH{} oferece: a transparência algorítmica como exigência ética, não apenas técnica.
\end{fichamento}

\begin{fichamento}
    \textbf{Ref.:} \cite{Edwards2017PC}
    
    \textbf{Conceitos-chave:} direito à explicação, GDPR, profiling, opacidade, autonomia, decisões automatizadas.
    
    \textbf{Conceitos principais:}
    \begin{itemize}
        \item A análise crítica do direito à explicação no GDPR: o alcance limitado do Art. 22 e sua baixa efetividade prática.
        \item A distinções entre diferentes tipos de decisões automatizadas (perfilamento, categorização, decisões individuais) e seus diferentes regimes de proteção.
        \item A opacidade dos sistemas de profiling como ameaça à autonomia pessoal — a pessoa é categorizada e julgada sem ter conhecimento ou controle sobre o processo.
        \item A proposta de um direito à intervenção humana significativa, não apenas formal, nas decisões automatizadas.
        \item A crítica à abordagem europeia de regulação da IA como insuficiente para proteger os direitos fundamentais dos cidadãos.
    \end{itemize}
    
    \textbf{Articulação:} A defesa do direito à explicação como garantia fundamental, aliada à crítica de suas limitações práticas no GDPR, fundamenta a análise da LGPD e dos projetos de lei brasileiros no capítulo quarto. A proposta de intervenção humana significativa dialoga com a exigência da \MH{} de que a técnica esteja sempre subordinada ao juízo ético da pessoa, jamais o substituindo de forma irreversível.
\end{fichamento}

\begin{fichamento}
    \textbf{Ref.:} \cite{Goodman2017EU}
    
    \textbf{Conceitos-chave:} regulação, tomada de decisão algorítmica, GDPR, direitos, transparência, não discriminação.
    
    \textbf{Conceitos principais:}
    \begin{itemize}
        \item A análise das disposições do GDPR relevantes para a tomada de decisão algorítmica, especialmente os Arts. 22, 13, 14 e 15.
        \item O direito à explicação como inovação do GDPR e seus fundamentos na tradição jurídica europeia de proteção de dados.
        \item A proibição de decisões exclusivamente automatizadas como regra geral, com exceções limitadas.
        \item As obrigações de transparência do controlador: informação significativa sobre a lógica do algoritmo.
        \item A avaliação crítica da eficácia do GDPR diante da complexidade técnica dos modelos de aprendizado de máquina.
    \end{itemize}
    
    \textbf{Articulação:} O artigo oferece a base comparativa europeia para a análise do regime jurídico brasileiro de regulação da IA no capítulo quarto. A estrutura de direitos e garantias do GDPR — especialmente o direito à explicação e a proibição de decisões exclusivamente automatizadas — constitui o parâmetro mínimo de proteção contra o qual a LGPD e o PL 2.338/2023 são avaliados, em diálogo com os princípios éticos propostos pela \MH{}.
\end{fichamento}

\begin{fichamento}
    \textbf{Ref.:} \cite{Binns2018ITG}
    
    \textbf{Conceitos-chave:} equidade, aprendizado de máquina, justiça algorítmica, imparcialidade, vieses, fairness.
    
    \textbf{Conceitos principais:}
    \begin{itemize}
        \item As diferentes definições formais de equidade (\textit{fairness}) no aprendizado de máquina: igualdade de oportunidades, paridade demográfica, imparcialidade individual.
        \item A incompatibilidade entre diferentes noções de equidade — não é possível satisfazer simultaneamente todos os critérios formais de justiça algorítmica.
        \item A necessidade de escolhas normativas explícitas sobre qual concepção de equidade orienta o desenvolvimento de cada sistema de IA.
        \item A tensão entre equidade algorítmica e acurácia preditiva — sistemas mais justos podem ser menos precisos e vice-versa.
        \item A importância do contexto social e institucional na determinação do que constitui uma decisão algorítmica justa.
    \end{itemize}
    
    \textbf{Articulação:} O artigo fornece critérios técnicos e conceituais para avaliar a equidade dos sistemas de IA, compatível com o princípio de justiça da \MH{}. A demonstração de que não existe uma única definição neutra de equidade algorítmica reforça a exigência da encíclica de que o desenvolvimento tecnológico seja orientado por uma visão integral da pessoa humana e do bem comum, e não por critérios meramente formais ou procedimentais.
\end{fichamento}

\begin{fichamento}
    \textbf{Ref.:} \cite{Selbst2018INC}
    
    \textbf{Conceitos-chave:} equidade, intuição, discricionariedade, justiça algorítmica, normas sociais, contexto.
    
    \textbf{Conceitos principais:}
    \begin{itemize}
        \item A tensão entre as noções intuitivas de equidade (\textit{intuitive notions of fairness}) e as definições formais de equidade algorítmica.
        \item A discricionariedade humana como elemento que sistemas automatizados eliminam, mas que pode ser necessária para decisões justas e contextualizadas.
        \item A diferença entre a flexibilidade interpretativa humana (que permite adaptar regras a casos específicos) e a rigidez dos sistemas algorítmicos.
        \item A perda de nuances contextuais na tradução de decisões sociais complexas para modelos matemáticos.
        \item A necessidade de preservar espaço para o juízo humano e a discricionariedade informada em decisões com impacto significativo sobre direitos.
    \end{itemize}
    
    \textbf{Articulação:} As limitações da equidade algorítmica identificadas por Selbst — especialmente a perda de discricionariedade contextual — informam a necessidade de uma abordagem normativa mais ampla como a proposta pela \MH{}. A crítica à rigidez dos sistemas automatizados reforça a defesa da encíclica pela irredutibilidade da pessoa a critérios algorítmicos e pela necessidade de preservar o juízo ético humano como instância última de decisão.
\end{fichamento}

\begin{fichamento}
    \textbf{Ref.:} \cite{Selbst2018TIR}
    
    \textbf{Conceitos-chave:} regulação, insuficiência, adequação, direitos, transparência, contexto.
    
    \textbf{Conceitos principais:}
    \begin{itemize}
        \item A inadequação dos marcos regulatórios existentes para lidar com os desafios específicos dos sistemas de IA.
        \item Os ``desalinhamentos'' (\textit{mismatches}) entre as características técnicas da IA e os pressupostos dos modelos regulatórios tradicionais.
        \item A identificação de cinco categorias de desalinhamento: tecnológico, sociotécnico, de agência, de conhecimento e de contexto.
        \item A necessidade de reformular as abordagens regulatórias para considerar a natureza sociotécnica dos sistemas algorítmicos.
        \item A proposta de uma regulação contextual e adaptativa, que reconheça as especificidades de cada domínio de aplicação da IA.
    \end{itemize}
    
    \textbf{Articulação:} A tese da inadequação dos marcos regulatórios existentes — em particular os cinco desalinhamentos identificados por Selbst — fundamenta a análise crítica do capítulo quarto sobre as insuficiências do ordenamento brasileiro diante dos desafios da IA. A proposta de regulação contextual dialoga com a abordagem da \MH{}, que também reconhece a necessidade de princípios éticos que orientem a aplicação da técnica em contextos específicos, sem cair em generalizações abstratas ou prescrições inflexíveis.
\end{fichamento}

\begin{fichamento}
    \textbf{Ref.:} \cite{Bostrom2005TH}
    
    \textbf{Conceitos-chave:} transumanismo, superinteligência, ética, futuro da humanidade, aperfeiçoamento humano, risco existencial.
    
    \textbf{Conceitos principais:}
    \begin{itemize}
        \item O transumanismo como movimento intelectual e cultural que defende o aperfeiçoamento da condição humana através da tecnologia.
        \item A superinteligência como objetivo último do transumanismo — uma inteligência artificial que supera o melhor cérebro humano em todos os domínios.
        \item Os riscos existenciais associados ao desenvolvimento descontrolado da IA — a possibilidade de extinção da humanidade.
        \item A crítica às objeções bioéticas e humanistas ao transumanismo, que Bostrom considera moralmente paralisantes.
        \item A defesa do aperfeiçoamento técnico do ser humano como extensão do projeto humanista de autossuperação.
    \end{itemize}
    
    \textbf{Articulação:} A análise do transumanismo por Bostrom — que defende o aperfeiçoamento humano através da tecnologia — informa a crítica da \MH{} ao paradigma tecnocrático. A defesa bostromiana da superação dos limites humanos pela técnica representa a visão antagônica à antropologia integral da encíclica, que afirma a irredutibilidade da pessoa à sua dimensão biológica ou informacional. O contraste entre as duas visões, desenvolvido no capítulo terceiro, permite demonstrar a originalidade e a atualidade da proposta ética da \MH{}.
\end{fichamento}

\section{Legislação e Documentos}

\begin{fichamento}
    \textbf{Ref.:} \cite{Brasil1988CF}
    
    \textbf{Conceitos-chave:} dignidade da pessoa humana, direitos fundamentais, personalidade, privacidade, honra, imagem, inviolabilidade.
    
    \textbf{Dispositivos relevantes:}
    \begin{itemize}
        \item Art. 1º, III: Dignidade da pessoa humana como fundamento da República.
        \item Art. 5º, X: Inviolabilidade da vida privada, honra e imagem.
        \item Art. 5º, LXXII: Habeas data como garantia fundamental.
    \end{itemize}
    
    \textbf{Articulação:} A Constituição Federal constitui o fundamento normativo máximo da proteção da personalidade, examinada no capítulo quarto em seus limites e possibilidades diante dos desafios da IA.
\end{fichamento}

\begin{fichamento}
    \textbf{Ref.:} \cite{Brasil2018LGPD}
    
    \textbf{Conceitos-chave:} proteção de dados, consentimento, direito à explicação, decisões automatizadas, ANPD, sanções.
    
    \textbf{Dispositivos relevantes:}
    \begin{itemize}
        \item Art. 6º: Princípios da proteção de dados (finalidade, adequação, necessidade, transparência, prevenção, não discriminação).
        \item Art. 20: Direito à revisão de decisões automatizadas.
        \item Art. 38: Obrigação de comunicação de incidentes de segurança.
    \end{itemize}
    
    \textbf{Articulação:} A LGPD é o principal instrumento jurídico brasileiro de proteção de dados pessoais. Sua análise crítica no capítulo quarto e a articulação com os princípios da \MH{} no capítulo quinto constituem o núcleo normativo da pesquisa.
\end{fichamento}

\begin{fichamento}
    \textbf{Ref.:} \cite{Brasil2023PL2338}
    
    \textbf{Conceitos-chave:} IA, regulação, classificação de risco, transparência, governança, accountability, fiscalização.
    
    \textbf{Dispositivos relevantes:}
    \begin{itemize}
        \item Classificação dos sistemas de IA por nível de risco.
        \item Obrigações de transparência e explicabilidade.
        \item Criação de autoridade competente para fiscalização.
    \end{itemize}
    
    \textbf{Articulação:} O PL 2.338/2023 representa a tentativa mais recente de regulação específica da IA no Brasil. Sua análise no capítulo quarto e a articulação com os princípios éticos da \MH{} no capítulo quinto apontam caminhos para seu aperfeiçoamento.
\end{fichamento}

\begin{fichamento}
    \textbf{Ref.:} \cite{Brasil2002CC}
    
    \textbf{Conceitos-chave:} pessoa natural, capacidade, direitos da personalidade, nome, imagem, honra, privacidade, disposição do próprio corpo.
    
    \textbf{Dispositivos relevantes:}
    \begin{itemize}
        \item Arts. 1º–10: Da pessoa natural — capacidade jurídica, nascimento, morte, ausência, domicílio.
        \item Arts. 11–21: Dos direitos da personalidade — intransmissibilidade, irrenunciabilidade, tutela do nome, imagem, honra e privacidade.
        \item Art. 12: Ação de cessação e indenização por violação a direitos da personalidade, inclusive após a morte.
        \item Art. 20: Proteção da imagem e da palavra — proibição de divulgação de escritos, palavras ou imagem sem autorização.
        \item Arts. 186–187: Responsabilidade civil por ato ilícito e abuso de direito como fundamento para reparação de danos à personalidade causados por sistemas de IA.
    \end{itemize}
    
    \textbf{Articulação:} O Código Civil fundamenta a proteção dos direitos da personalidade no direito civil brasileiro, examinada no capítulo quarto. A abertura do sistema de direitos da personalidade — que a doutrina reconhece como rol exemplificativo — permite a tutela de novas manifestações da personalidade ameaçadas pela IA, como a autodeterminação informativa e o direito à explicação algorítmica, em harmonia com a visão integral da pessoa proposta pela \MH{}.
\end{fichamento}

\begin{fichamento}
    \textbf{Ref.:} \cite{UE2016GDPR}
    
    \textbf{Conceitos-chave:} proteção de dados, privacidade, GDPR, decisões automatizadas, direito à explicação, transferência internacional, consentimento.
    
    \textbf{Dispositivos relevantes:}
    \begin{itemize}
        \item Arts. 13–14: Obrigações de transparência do controlador — informações sobre a lógica do tratamento e as consequências previsíveis para o titular.
        \item Art. 22: Direito do titular a não ficar sujeito a decisões exclusivamente automatizadas, incluindo perfilamento, com exceções limitadas e garantias.
        \item Arts. 35–36: Obrigação de realizar avaliação de impacto sobre a proteção de dados (\textit{Data Protection Impact Assessment}) para tratamentos de alto risco.
        \item Art. 25: Proteção de dados desde a concepção e por padrão (\textit{privacy by design and by default}).
        \item Art. 46: Transferência internacional de dados pessoais — garantias adequadas e decisões de adequação.
    \end{itemize}
    
    \textbf{Articulação:} O GDPR é o marco regulatório de referência internacional para proteção de dados, que influencia a LGPD e o PL 2.338/2023 no Brasil. Sua análise comparativa no capítulo quarto permite identificar tanto os avanços da legislação brasileira quanto suas lacunas face ao direito europeu, enquanto os princípios de transparência, privacidade desde a concepção e avaliação de impacto dialogam com a exigência da \MH{} de subordinar o desenvolvimento técnico ao controle ético e à proteção integral da pessoa humana.
\end{fichamento}

\begin{fichamento}
    \textbf{Ref.:} \cite{STF2020ADI6640}
    
    \textbf{Conceitos-chave:} proteção de dados, direito fundamental, ADI 6640, STF, LGPD, dignidade, privacidade.
    
    \textbf{Dispositivos relevantes:}
    \begin{itemize}
        \item Voto do relator, Ministro Luiz Fux: A proteção de dados pessoais como direito fundamental autônomo, extraído do princípio da dignidade da pessoa humana e do direito à privacidade.
        \item Tese firmada: A LGPD é constitucional e a proteção de dados é direito fundamental com eficácia imediata, nos termos do Art. 5º, §2º da CF.
        \item Fundamento: A autodeterminação informativa como projeção do livre desenvolvimento da personalidade e da autonomia privada.
        \item Reconhecimento da competência concorrente da União para legislar sobre proteção de dados, superando controvérsias federativas.
        \item Precedente de modulação: O STF estabeleceu o prazo de 24 meses para adequação à LGPD, reconhecendo a complexidade da transição.
    \end{itemize}
    
    \textbf{Articulação:} O STF, na ADI 6640, reconheceu a proteção de dados como direito fundamental autônomo, fundamentando a necessidade de regulação da IA no Brasil. O voto do relator, que conecta a proteção de dados à dignidade da pessoa humana e ao livre desenvolvimento da personalidade, dialoga diretamente com a fundamentação antropológica da \MH{}, que também vincula a proteção da pessoa na era digital à sua dignidade intrínseca e irredutível.
\end{fichamento}

\begin{fichamento}
    \textbf{Ref.:} \cite{STF2022ADPF695}
    
    \textbf{Conceitos-chave:} compartilhamento de dados, administração pública, privacidade, proteção de dados, limites ao poder estatal, LGPD.
    
    \textbf{Dispositivos relevantes:}
    \begin{itemize}
        \item Decisão do STF que estabeleceu limites ao compartilhamento de dados entre órgãos da administração pública federal.
        \item Exigência de base legal específica e de transparência para o compartilhamento de dados pessoais entre entidades públicas.
        \item Reafirmação da aplicabilidade da LGPD ao setor público, inclusive no que tange às sanções por violação.
        \item Necessidade de avaliação de impacto à proteção de dados antes de implementar sistemas de compartilhamento em larga escala.
        \item Vinculação do tratamento de dados pela administração pública aos princípios constitucionais da impessoalidade, moralidade e eficiência.
    \end{itemize}
    
    \textbf{Articulação:} A decisão do STF na ADPF 695 reafirmou limites à utilização de dados pela administração pública, tema relevante para a regulação da IA no setor público, examinada no capítulo quarto. O reconhecimento de que o poder público está vinculado à LGPD e aos princípios constitucionais quando trata dados pessoais dialoga com a exigência da \MH{} de que o desenvolvimento tecnológico seja orientado pelo bem comum e pelo respeito à dignidade de cada pessoa, especialmente aquelas em situação de vulnerabilidade.
\end{fichamento}

\begin{fichamento}
    \textbf{Ref.:} \cite{HighLevelExpertGroup2019EA}
    
    \textbf{Conceitos-chave:} IA confiável, licitude, eticidade, robustez técnica, transparência, accountability, bem-estar social.
    
    \textbf{Dispositivos relevantes:}
    \begin{itemize}
        \item Cap. I: Requisitos para uma IA confiável — ação humana e supervisão, robustez técnica e segurança, privacidade e governança de dados, transparência, diversidade e não discriminação, bem-estar social e ambiental, accountability.
        \item Cap. II: Avaliação de confiabilidade — lista de verificação (\textit{assessment list}) para desenvolvedores e implementadores de sistemas de IA.
        \item Cap. III: Governança da IA — necessidade de mecanismos de governança multinível que articulem autorregulação, regulação estatal e participação social.
        \item Recomendações para políticas públicas: investimento em pesquisa, educação digital, fortalecimento de autoridades reguladoras e cooperação internacional.
        \item Ênfase na centralidade da pessoa (\textit{human-centric AI}) como princípio orientador de todo desenvolvimento tecnológico.
    \end{itemize}
    
    \textbf{Articulação:} As diretrizes éticas do Grupo de Peritos de Alto Nível da UE constituem o principal documento internacional de ética da IA, analisado comparativamente com a \MH{} no capítulo segundo. A centralidade da pessoa (\textit{human-centric AI}) como princípio orientador converge com a visão antropológica da encíclica, enquanto os sete requisitos para IA confiável oferecem um quadro operacional que pode ser enriquecido pelos princípios éticos mais amplos propostos pela \MH{}, especialmente o princípio da irredutibilidade da pessoa e o destino universal dos bens aplicado aos dados.
\end{fichamento}


\end{anexosenv}

\end{document}
